% 木星（综述）
% license CCBYSA3
% type Wiki

本文根据 CC-BY-SA 协议转载翻译自维基百科\href{https://en.wikipedia.org/wiki/Jupiter}{相关文章}。

\begin{figure}[ht]
\centering
\includegraphics[width=8cm]{./figures/572db8cd95497061.png}
\caption{木星的真实色彩图像\(^\text{[a]}\)，由哈勃空间望远镜于 2024 年 1 月拍摄\(^\text{[b]}\)（Jupiter in true colour\(^\text{[a]}\), taken by the Hubble Space Telescope in January 2024\(^\text{[b]}\)）} \label{fig_Jupite_20}
\end{figure}
木星是距离太阳第五颗行星，也是太阳系中体积最大的行星。它是一颗气态巨行星，其质量几乎是太阳系中所有其他行星总和的 2.5 倍，但仍不到太阳质量的千分之一。它的直径是地球的 11 倍，是太阳的十分之一。木星以 5.20 天文单位（778.5 吉米）的平均距离绕太阳运行，公转周期为 11.86 年。在地球夜空中，木星是仅次于月球和金星的第三亮的天然天体，自史前时代以来便为人类所观察。它的名字源自古罗马宗教中至高神朱庇特（Jupiter）。

木星是太阳行星中最先形成的一颗，其在太阳系原始阶段向内迁移的过程深刻影响了其他行星的形成历史。木星的大气按质量计由 76\% 的氢和 24\% 的氦组成，内部密度更高。它还含有微量的碳、氧、硫、氖等元素，以及氨、水蒸气、磷化氢、硫化氢和烃类等化合物。木星的氦含量约为太阳的 80\%，与土星的成分相似。

木星外层大气被分为一系列纬向云带，这些云带在交界处产生湍流和风暴；其中最显著的结果就是“大红斑”——自 1831 年以来被记录的一场巨型风暴。由于木星自转极快（约 10 小时自转一周），使其呈现为一颗扁球体；与两极相比，赤道部位存在约 6.5\% \(^\text{[e]}\)的明显隆起。科学家认为，木星内部结构包括一层由流体金属氢组成的外地幔，以及由更高密度物质构成的弥散内核。木星内部持续的收缩过程释放出比其从太阳接收的还要多的热量。木星的磁场是太阳系中最强且连续范围第二大的结构，由流体金属氢内部的涡电流产生。太阳风与其磁层相互作用，使磁层向外延伸并影响木星的运行环境。

至少有 97 颗卫星绕木星运行；其中四颗最大的卫星——木卫一（Io）、木卫二（Europa）、木卫三（Ganymede）和木卫四（Callisto）——位于其磁层内，用普通双筒望远镜即可观测。四者中最大的木卫三比行星水星还要大。木星周围还存在一个暗淡的行星环系统。木星环主要由尘埃组成，分为三部分：内侧呈光环状尘粒环的“晕环”（halo）、相对明亮的主环（main ring），以及外侧的薄雾状“薄环”（gossamer ring）。这些行星环在可见光和近红外光下呈现红色。木星环系统的年龄尚不清楚，可能追溯至木星形成时期。自 1973 年以来，已有九个机器人探测器访问过木星：包括七次飞掠任务和两艘专门的轨道探测器（另有两艘在途中）。在其他行星系统中也已发现类似木星的系外行星。
\subsection{名称与符号}.
Jupiter 在古希腊和古罗马文明中都以神系中的最高神命名：对希腊人而言是宙斯（Zeus），对罗马人而言是朱庇特（Jupiter）\(^\text{[20]}\)。国际天文学联合会于 1976 年正式采用 Jupiter 作为该行星的名称，并自此将新发现的卫星以这位神祇的情人、宠爱者及后代命名\(^\text{[21]}\)。木星的行星符号 ♃ 源自带有水平笔画的希腊字母 zeta（⟨Ƶ⟩），作为 Zeus 的缩写\(^\text{[22][23]}\)。在拉丁语中，Iovis 是 Iuppiter（即 Jupiter）的属格形式，并与 Zeus（“天空之父”）的词源相关。英语中的对应形式 Jove 被认为在 14 世纪起作为该行星的诗性名称出现\(^\text{[24]}\)。Jovian 是 Jupiter 的形容词形式。较早的形容词 jovial，被中古时代的占星家使用，后来演变为“愉快的”或“欢乐的”含义，此情绪被认为与木星在占星学中的影响相关\(^\text{[25]}\)。原始希腊神祇 Zeus 还提供了词根 zeno-，用于构成一些与木星相关的词语，例如 zenography\(^\text{[f]}\)。
\subsection{形成与迁移}
木星被认为是太阳系中最古老的行星，其形成时间仅在太阳之后一百万年，约早于地球五千万年\(^\text{[26]}\)。当前的太阳系形成模型表明，木星形成于雪线处或其以外的位置：即距离早期太阳足够远、温度足够低，使水等挥发物能够凝结成固体的区域\(^\text{[27]}\)。木星首先形成固体核心，随后吸积其气态大气层。因此，该行星必须在太阳星云完全消散之前形成\(^\text{[28]}\)。在形成过程中，木星的质量逐渐增加，直到达到地球质量的 20 倍，其中约一半由硅酸盐、冰以及其他重元素成分构成\(^\text{[26]}\)。当原始木星增大超过 50 个地球质量时，它在太阳星云中开辟出一个空隙\(^\text{[26]}\)。此后，这颗不断成长的行星在 300–400 万年内达到最终质量\(^\text{[26][28]}\)。

根据“Grand tack 假说”，木星最初在距太阳约 3.5 AU（5.20 亿公里；3.30 亿英里）的位置开始形成。随着年轻的木星吸积质量，它与环绕太阳的气体盘相互作用，并受到土星轨道共振的影响，从而向内迁移\(^\text{[27][29]}\)。这扰乱了数颗在更靠近太阳轨道上运行的超级地球，使它们发生破坏性的碰撞\(^\text{[30]}\)。随后，土星开始以比木星更快的速度向内迁移，直到两者在距太阳约 1.5 AU（2.20 亿公里；1.40 亿英里）处被俘获进入 3:2 平均运动共振\(^\text{[31]}\)。这一事件改变了迁移方向，使它们开始远离太阳并迁出内太阳系，移动至如今的位置\(^\text{[30]}\)。上述全过程持续约 300–600 万年，而木星最终阶段的迁移历经数十万年完成\(^\text{[29][32]}\)。木星从内太阳系的迁移最终使包括地球在内的内行星得以从残余碎片中形成\(^\text{[33]}\)。

Grand tack 假说仍存在多个未解决的问题。其推导出的类地行星形成时间尺度似乎与测量得到的元素组成不一致\(^\text{[34]}\)。如果木星曾穿越太阳星云迁移，它最终可能会定居在更靠近太阳的轨道上\(^\text{[35]}\)。一些太阳系形成的竞争模型预测木星的形成轨道性质与今日行星的位置非常接近\(^\text{[28]}\)。其他模型则预测木星形成于更为外侧的位置，例如 18 AU（27 亿公里；17 亿英里）\(^\text{[36][37]}\)。

根据 Nice 模型，在太阳系历史最初的 6 亿年间，原始柯伊伯带天体的内落导致木星和土星从它们最初的位置迁移至 1:2 共振，这使土星移入更高轨道，从而扰乱了天王星和海王星的轨道，消耗了柯伊伯带，并触发了晚期重轰炸\(^\text{[38]}\)。

根据 Jumping-Jupiter 情景，木星在早期太阳系中的迁移可能导致第五颗气态巨行星被抛出。该假说认为，在其轨道迁移过程中，木星的引力扰动破坏了其他气态巨行星的轨道，可能使其中一颗完全被逐出太阳系。此类事件的动力学将极大改变太阳系的形成与配置，最终只留下如今人类所见的四颗气态巨行星\(^\text{[39]}\)。

基于木星的组成，有研究者提出其最初形成于分子氮（N\(_2\)）雪线之外，即距太阳 20–30 AU（30–45 亿公里；19–28 亿英里）的位置，甚至可能形成于氩雪线之外，后者可能远至 40 AU（60 亿公里；37 亿英里）\(^\text{[40][41]}\)。如果木星形成于这些极远的位置之一，则它将在大约 70 万年的时间尺度上向内迁移至目前的位置\(^\text{[36][37]}\)，这一过程发生于行星开始形成后约 200–300 万年的时代。在该模型中，土星、天王星和海王星的形成位置将比木星更远，而土星亦会向内迁移\(^\text{[36]}\)。
\subsection{物理特征}
木星是一颗气态巨行星，这意味着其化学组成主要为氢和氦。这些材料在行星地质学中被分类为“气体”，这一术语并不指示其物质状态。木星是太阳系中最大的行星，其赤道直径为 142,984 公里（88,846 英里），体积为地球的 1,321 倍\(^\text{[3][42]}\)。其平均密度为 1.326 g/cm\(^3\)\(^\text{[g]}\)，低于四颗类地行星的密度\(^\text{[44][45]}\)。
\subsubsection{成分}
木星大气按质量计约由 76\% 的氢和 24\% 的氦组成。按体积计，上层大气约含 90\% 的氢和 10\% 的氦，这一较低比例是因为单个氦原子比该大气层中形成的氢分子更为致密\(^\text{[46]}\)。大气含有微量的碳、氧、硫和氖\(^\text{[47]}\)，以及氨、水蒸气、磷化氢、硫化氢和甲烷、乙烷、苯等烃类化合物\(^\text{[48]}\)。其最外层包含冻结氨的晶体\(^\text{[49]}\)。行星内部密度更高，按质量计约由 71\% 的氢、24\% 的氦和 5\% 的其他元素组成\(^\text{[50][51]}\)。

大气中氢与氦的比例接近理论上的原始太阳星云成分\(^\text{[52]}\)。上层大气中的氖含量为每百万份质量中的 20 份，约为太阳丰度的十分之一\(^\text{[53]}\)。木星的氦含量约为太阳的 80\%，这是由于这些元素在行星内部深处以富含氦的液滴形式沉降的过程所致\(^\text{[54][55]}\)。

基于光谱学分析，土星的组成被认为与木星相似，但另外两颗巨行星天王星和海王星含有相对更少的氢和氦，而相对含有更多的次常见元素，包括氧、碳、氮和硫\(^\text{[56]}\)。这些行星被称为冰巨行星，因为在其形成期间，这些元素被认为以冰的形式并入其中；然而它们内部实际上可能含有极少量的冰\(^\text{[57]}\)。
\subsubsection{大小与质量}
\begin{figure}[ht]
\centering
\includegraphics[width=8cm]{./figures/82e5bc11e94dd9a1.png}
\caption{木星与地球及地球月球的尺寸比较} \label{fig_Jupite_1}
\end{figure}
木星的直径约为地球的十一倍（11.209 R🜨）；其质量是地球的 318 倍\(^\text{[3]}\)，相当于太阳系中其他所有行星质量总和的 2.5 倍。其质量之大，以至于它与太阳的质心位于太阳表面之外，距太阳中心 1.068 个太阳半径\(^\text{[58][59]}\)。木星的半径约为太阳半径的十分之一（0.10276 R☉）\(^\text{[60]}\)，其质量约为太阳质量的千分之一，两者的平均密度相似\(^\text{[61]}\)。“木星质量”（MJ 或 MJup）常被用作描述其他天体质量的单位，特别是系外行星和棕矮星。例如，系外行星 HD 209458 b 的质量为 0.69 MJ，而棕矮星 Gliese 229 b 的质量为 60.4 MJ\(^\text{[62][63]}\)。

木星辐射的热量多于其从太阳接收的辐射能量，这是因为其内部收缩过程中的 Kelvin–Helmholtz 机制\(^\text{[64]:30^\text{}[65]}\)。这一过程导致木星每年缩小约 1 毫米（0.039 英寸）\(^\text{[66][67]}\)。在其形成之初，木星更为炽热，且其直径约为当前的两倍\(^\text{[68][69]}\)。

理论模型表明，如果木星再增加 40\% 以上的质量，其内部将因压缩而导致体积缩小，尽管质量增加。对较小的质量变化而言，其半径变化并不显著\(^\text{[70]}\)。因此，木星被认为已达到由其组成与演化历史所能达到的最大行星直径\(^\text{[71]}\)。随着质量进一步增加而导致的收缩过程会持续，直至达到明显的恒星点火为止\(^\text{[72]}\)。尽管木星需要约 75 倍目前质量才能点燃氢聚变、成为一颗恒星\(^\text{[73]}\)，但其直径依然符合最小红矮星可能仅略大于土星半径的事实\(^\text{[74]}\)。
\subsubsection{内部结构}
\begin{figure}[ht]
\centering
\includegraphics[width=8cm]{./figures/2dfe73e2dbd3a048.png}
\caption{木星的内部结构、表面特征、行星环和内侧卫星示意图} \label{fig_Jupite_2}
\end{figure}
在 21 世纪初之前，大多数科学家提出了两种关于木星形成的主要情景之一。如果行星首先作为固体天体吸积形成，它将由一个致密核心构成，其外包围一层流体金属氢（含少量氦），这一层向外延伸至行星半径的约 80\%\(^\text{[75]}\)，并有一层主要由分子氢组成的外层大气\(^\text{[67]}\)。另一种可能是，如果行星直接从气态原行星盘中坍缩形成，则预计它将完全缺乏核心，而是从中心向外由逐渐变得更致密的流体构成（主要是分子氢与金属氢）。Juno 任务的数据表明，木星具有一个向地幔混合的弥散核心，延伸至行星半径的 30–50\%，由总质量相当于 7–25 个地球质量的重元素构成\(^\text{[76][77][78]}\)。这种混合过程可能在形成期间便已发生，当时行星从周围星云中吸积固体与气体\(^\text{[79]}\)。另一种可能是，在木星形成几百万年后，一颗约十倍地球质量的行星发生碰撞，扰乱了原本致密的木星核心\(^\text{[77][80]}\)。

在金属氢层之外，是透明的内部氢大气。在这一深度，压力和温度高于分子氢的临界压力 1.3 MPa 与临界温度 33 K（−240.2 °C；−400.3 °F）\(^\text{[81]}\)。在这种状态下，不再存在明确的液态与气态的分界——氢处于所谓的超临界流体状态。从云层向下延伸的氢气和氦气在更深处逐渐过渡为液体，可能类似于由液态氢及其他超临界流体构成的“海洋”\(^\text{[64]:22^\text{}[82][83]}\)。从物理上讲，随着深度增加，气体会逐渐变得更热、更致密\(^\text{[84][85]}\)。

氦和氖如同“雨滴”般在大气下层向下沉降，使上层大气中的这些元素含量减少\(^\text{[54][86]}\)。计算表明，氦滴在距中心 60,000 公里（37,000 英里）（即云顶下方 11,000 公里［6,800 英里］）处从金属氢中分离，并在 50,000 公里（31,000 英里）（云层下方 22,000 公里［14,000 英里］）处重新混合\(^\text{[87]}\)。研究指出，如同土星\(^\text{[88]}\)以及冰巨行星天王星和海王星\(^\text{[89]}\)一样，木星内部也可能存在钻石雨。

木星内部的温度和压力向内稳步升高，因为行星形成时的热量只能通过对流方式散失\(^\text{[55]}\)。在大气压力约为一个地球标准大气压（约 0.10 MPa（1 bar））的深度，温度约为 165 K（−108 °C；−163 °F）。超临界氢从分子流体逐渐过渡到金属流体的区域涵盖 50–400 GPa 的压力范围，对应的温度为 5,000–8,400 K（4,730–8,130 °C；8,540–14,660 °F）。木星弥散核心的温度估计为 20,000 K（19,700 °C；35,500 °F），压力约为 4,000 GPa\(^\text{[90]}\)。
\subsubsection{大气}
木星大气主要由分子氢和氦组成，含有较少量的水、甲烷、硫化氢和氨等其他化合物\(^\text{[91]}\)。木星大气自云层向下延伸约 3,000 公里（2,000 英里）\(^\text{[90]}\)。

\textbf{云层}
\begin{figure}[ht]
\centering
\includegraphics[width=8cm]{./figures/a270704535c387fb.png}
\caption{木星云层系统在一个月中的运动延时影像（1979 年旅行者 1 号飞掠期间拍摄）} \label{fig_Jupite_3}
\end{figure}
木星上始终覆盖着由氨晶体构成的云层，其中可能还含有硫氢化铵\(^\text{[92]}\)。这些云位于大气的对流层顶，形成分布在不同行星纬度上的云带，被称为热带区域。它们进一步被划分为较浅色的“区”（zones）与较深色的“带”（belts）。这些相互冲突的环流模式的相互作用引发风暴与湍流。在纬向急流中，常见的风速可达每秒 100 米（360 公里/小时；220 英里/小时）\(^\text{[93]}\)。这些云带的宽度、颜色与强度会因年份而变化，但其整体结构足够稳定，使科学家能够为它们命名\(^\text{[59]:6^\text{}]}\)。

云层厚约 50 公里（31 英里），由至少两层氨云构成：顶部为一层较薄、较透明的区域，下方为更厚的下层云带。可能在氨云之下还存在一薄层水云，这一推测来自对木星大气中闪电现象的探测\(^\text{[94]}\)。这些放电的能量可高达地球闪电的千倍\(^\text{[95]}\)。水云被认为以与地球雷暴相似的方式形成雷暴，由木星内部上升的热量驱动\(^\text{[96]}\)。Juno 任务揭示了“浅层闪电”的存在，其起源于大气中较高高度的氨-水云\(^\text{[97]}\)。这些放电会携带由水-氨泥浆构成、外覆冰壳的“浆球”（mushballs），它们会坠入大气深处\(^\text{[98]}\)。在木星高层大气中还观测到类似地球“精灵”与“红色妖精”的上层大气闪电——持续约 1.4 毫秒的亮光，由于氢元素的缘故呈现蓝色或粉红色\(^\text{[99][100]}\)。

木星云层中的橙色与棕色由下沉的化合物造成，这些化合物在受到来自太阳的紫外光照射后发生变色。其具体成分尚不确定，但普遍认为由磷、硫或可能的烃类构成\(^\text{[64]:39^\text{}[101]}\)。这些被称为“着色团”的化合物与来自下层的暖云混合。浅色云区则是因上升对流产生氨晶体，这些晶体将着色团遮蔽于视线之外\(^\text{[102]}\)。

木星自转轴倾角很小，因此其两极永远接收的太阳辐射比赤道区域少。行星内部的对流将能量输送至两极，从而在云层高度维持热量平衡\(^\text{[59]:54^\text{}]}\)。

\textbf{大红斑及其他涡旋}
\begin{figure}[ht]
\centering
\includegraphics[width=8cm]{./figures/b90e5abcdffbaa73.png}
\caption{由朱诺号探测器以真实色彩拍摄的大红斑特写。由于朱诺号的成像方式，拼接后的图像呈现出显著的桶形畸变。} \label{fig_Jupite_4}
\end{figure}
木星最著名的特征之一是大红斑\(^\text{[103]}\)，这是一场位于南纬 22° 的持久反气旋风暴。它首次于 1831 年被观测到\(^\text{[104]}\)，甚至可能早在 1665 年就已被记录\(^\text{[105][106]}\)。哈勃太空望远镜所获得的图像显示，在大红斑附近还有另外两个“红斑”\(^\text{[107][108]}\)。通过口径 12 厘米或更大的地基望远镜即可观测到这场风暴\(^\text{[109]}\)。风暴呈逆时针方向旋转，其周期约为六天\(^\text{[110]}\)。该风暴的最高高度约比周围云顶高出 8 公里（5 英里）\(^\text{[111]}\)。大红斑的具体成分及其红色来源仍不确定，不过光解氨与乙炔发生反应是可能的解释之一\(^\text{[112]}\)。

大红斑的面积大于地球\(^\text{[113]}\)。数学模型表明，这场风暴稳定存在，将成为木星上的长期特征\(^\text{[114]}\)。然而，从发现至今它的尺寸已显著缩小。19 世纪末的早期观测显示其宽度约为 41,000 公里（25,500 英里）。截至 2015 年，风暴的大小约为 16,500 × 10,940 公里（10,250 × 6,800 英里）\(^\text{[115]}\)，且其长度每年缩减约 930 公里（580 英里）\(^\text{[113]}\)。2021 年 10 月，Juno 飞掠任务测量了大红斑的深度，结果约为 300–500 公里（190–310 英里）\(^\text{[116]}\)。

Juno 任务在木星两极发现了数个气旋群。北极气旋群包含九个气旋，其中一个位于中心，周围分布着另外八个；南极气旋群同样有一个中心涡旋，但周围由五个大型风暴与一个较小的风暴组成，总计七个\(^\text{[117][118]}\)。

在 2000 年，木星南半球出现了一个类似大红斑但规模较小的大气特征。这是由几个较小的白色椭圆风暴合并而成——这三个较小白色椭圆最初形成于 1939–1940 年。合并后的结构被命名为 Oval BA。此后，它的强度有所增强，颜色也由白色转为红色，因此获得了“小红斑”之称\(^\text{[119][120]}\)。

2017 年 4 月，在木星北极的热层中发现了一个“大冷斑”。这一结构宽 24,000 公里（15,000 英里），高 12,000 公里（7,500 英里），其温度比周围材料低 200 °C（360 °F）。虽然该冷斑在短期内形态与强度会变化，但其在大气中的总体位置已保持超过 15 年。它可能是一种类似大红斑的巨大涡旋，并似乎像地球热层中的涡旋一样呈准稳定状态。该结构可能由来自木卫一的带电粒子与木星强磁场相互作用形成，导致热量流动被重新分配\(^\text{[121]}\)。
\subsubsection{磁层}
\begin{figure}[ht]
\centering
\includegraphics[width=8cm]{./figures/3a2e5013a5750a65.png}
\caption{伽利略卫星对木星磁层的影响} \label{fig_Jupite_5}
\end{figure}
木星的磁场是太阳系所有行星中最强的\(^\text{[102]}\)，其偶极矩为 4.170 高斯（0.4170 mT），并相对于自转轴倾斜 10.31°。其表面磁场强度从 2 高斯（0.20 mT）到 20 高斯（2.0 mT）不等\(^\text{[122]}\)。该磁场被认为由流体金属氢核心内部的涡电流——导电物质的旋转运动——所产生。在距离木星约 75 个木星半径处，磁层与太阳风的相互作用产生弓形激波。环绕木星磁层外侧的是磁暂停层，其位置位于磁鞘层的内缘——这是位于磁暂停层与弓形激波之间的区域。太阳风与这些区域相互作用，使磁层在木星背风侧被拉长并向外延伸，几乎抵达土星轨道。木星最大的四颗卫星全部在磁层内运行，从而受到保护免受太阳风影响\(^\text{[64]:69^\text{}]}\)。

木卫一的火山喷发大量释放二氧化硫，在其轨道附近形成一个气体环。气体在木星磁层中被电离，产生硫和氧离子。它们与来自木星大气的氢离子一起，在木星赤道平面形成一个等离子体片。该等离子体片随着行星共同旋转，使偶极磁场被拉伸为磁盘状。等离子体片中的电子产生强烈的无线电信号，其中叠加着 0.6–30 MHz 范围内的短促爆发，这些信号可通过地面消费级短波无线电接收机探测到\(^\text{[123][124]}\)。当木卫一穿过该气体环时，其相互作用会产生阿尔芬波，将电离物质带入木星极区。因此，通过回旋脉塞机制产生无线电波，能量沿锥形表面传播。当地球穿过该锥面时，来自木星的无线电辐射可超过太阳的无线电输出\(^\text{[125]}\)。
\subsubsection{行星环}
\begin{figure}[ht]
\centering
\includegraphics[width=8cm]{./figures/143fda39a5f08d29.png}
\caption{以红外光拍摄的木星图像，显示其微弱的行星环，以及两颗卫星——Amalthea 与 Adrastea、极光及大气特征。} \label{fig_Jupite_6}
\end{figure}
木星拥有一个由三部分组成的微弱行星环系统：最内侧由颗粒构成的光环状“晕环”（halo），中间较为明亮的主环，以及最外侧的薄雾环（gossamer ring）\(^\text{[126]}\)。这些行星环似乎由尘埃组成，而土星的行星环则由冰构成\(^\text{[64]:65^\text{}]}\)。主环极可能由卫星 Adrastrea 与 Metis 被撞击后喷射出的物质构成，这些物质因木星强大的引力影响而被吸引进入行星环中。新的物质则通过额外撞击持续补充\(^\text{[127]}\)。类似地，Thebe 与 Amalthea 两颗卫星被认为产生了薄雾环中的两个不同尘埃成分\(^\text{[127]}\)。有证据显示存在第四条行星环，可能由 Amalthea 碰撞产生的碎片构成，并沿着该卫星的轨道分布\(^\text{[128]}\)。
\subsection{轨道与自转}
\begin{figure}[ht]
\centering
\includegraphics[width=8cm]{./figures/9931f8d9015015fb.png}
\caption{展示木星自转及其卫星轨道运动的三小时延时影像} \label{fig_Jupite_7}
\end{figure}
木星是唯一一颗其与太阳的质心位于太阳体积之外的行星，不过仅超过太阳半径的 7\%\(^\text{[129][130]}\)。木星与太阳之间的平均距离为 7.78 亿公里（5.20 AU），其绕太阳公转一周需 11.86 年。这大约是土星轨道周期的五分之二，构成一个近似的轨道共振\(^\text{[131]}\)。木星轨道平面相对于地球倾斜 1.30°。由于其轨道偏心率为 0.049，木星在近日点时比在远日点时更接近太阳约 7,500 万公里\(^\text{[3]}\)，意味着其轨道几乎是圆形的。与此相比，系外行星的发现揭示许多木星大小的行星具有高轨道偏心率。模型表明，这可能是因为太阳系中存在两颗巨行星；而拥有三颗或更多巨行星的系统更倾向于产生较大的偏心率\(^\text{[132]}\)。

木星的自转轴倾角为 3.13°，相对较小，因此其季节变化远不及地球与火星显著\(^\text{[133]}\)。

木星的自转速度是太阳系所有行星中最快的，自转周期略短于十小时；这导致其赤道隆起，通过业余望远镜便可轻松观测。由于木星并非固体，其上层大气呈现差异自转。木星极区大气的自转速度比赤道大气慢约五分钟\(^\text{[134]}\)。木星是一颗扁球体，这意味着其赤道直径比两极之间的直径更长\(^\text{[85]}\)。在木星上，赤道直径比极径长 9,276 公里（5,764 英里）\(^\text{[3]}\)。

跟踪行星自转时，尤其是在绘制大气特征运动时，会使用三个参考系统。系统 I 适用于 7°N 至 7°S 之间的纬度，其自转周期为行星最短，为 9 小时 50 分 30.0 秒。系统 II 适用于其余纬度，自转周期为 9 小时 55 分 40.6 秒\(^\text{[135]}\)。系统 III 由射电天文学家定义，对应于木星磁层的自转；其周期被视为木星的官方自转周期\(^\text{[136]}\)。
\subsection{观测}
\begin{figure}[ht]
\centering
\includegraphics[width=8cm]{./figures/a02a6f8994823765.png}
\caption{通过业余望远镜观测到的木星及其四颗伽利略卫星} \label{fig_Jupite_8}
\end{figure}
木星通常是夜空中第四亮的天体（仅次于太阳、月亮和金星）\(^\text{[102]}\)，尽管在冲日附近火星有时会比木星更亮。根据木星相对于地球的位置不同，其视星等可在冲日时亮至 −2.94，而在与太阳合相时暗至 −1.66\(^\text{[17]}\)。其平均视星等为 −2.20，标准差为 0.33\(^\text{[17]}\)。木星的视直径也随之变化，在 50.1 至 30.5 角秒之间\(^\text{[3]}\)。当木星经过轨道近日点时会出现有利的冲日，使其更接近地球\(^\text{[137]}\)。在冲日前后，木星会出现大约 121 天的逆行运动，在天空中逆向移动约 9.9°，之后恢复顺行\(^\text{[138]}\)。

由于木星的轨道位于地球轨道之外，从地球观测到的木星相位角始终小于 11.5°；因此，通过地基望远镜看到的木星总是几乎完全被照亮的。只有在飞掠木星的航天器任务中，才获得了呈弦月状的木星影像\(^\text{[139]}\)。在小型望远镜下通常可以看到木星的四颗伽利略卫星，以及横跨木星大气的云带。口径约 4–6 英寸（10–15 厘米）的较大望远镜在大红斑朝向地球时能够观测到该结构\(^\text{[140][141]}\)。
\subsubsection{历史}
\textbf{望远镜出现前的研究}
\begin{figure}[ht]
\centering
\includegraphics[width=8cm]{./figures/b20ffb6befc24739.png}
\caption{《至大论》（Almagest）中木星（☉）相对于地球（🜨）经度运动的模型} \label{fig_Jupite_9}
\end{figure}
对木星的观测在公元前 7–8 世纪的巴比伦天文学家记录中已有确证\(^\text{[142]}\)。古代中国人称木星为“岁星”（Suìxīng 歲星），并依据木星绕太阳一周约需十二年左右的时间建立了十二地支体系；汉语至今仍以“岁”字指代年龄的“年”。到了公元前 4 世纪，这些观测发展为中国的黄道十二宫\(^\text{[143]}\)，每一年都与一位“太岁”星和掌管木星在夜空位置相对区域的天神相关联。这些信仰在部分道教与民间宗教实践以及东亚黄道的十二生肖中仍有保留。中国历史学家席泽宗提出，古代天文学家甘德\(^\text{[144]}\)曾报告过一颗与木星“相伴”的小星\(^\text{[145]}\)，这可能意味着他曾以肉眼观测到木星的一颗卫星。如果属实，这一发现将比伽利略的观测早近两千年\(^\text{[146][147]}\)。

2016 年的一篇论文指出，巴比伦人在公元前 50 年以前便使用梯形法则来对木星在黄道上的速度进行积分\(^\text{[148]}\)。在公元 2 世纪的著作《至大论》（Almagest）中，希腊化时期的天文学家托勒密（Claudius Ptolemaeus）构建了一个基于本轮与均轮的地心行星模型，用以解释木星相对于地球的运动，并给出其绕地球公转周期为 4332.38 天，即 11.86 年\(^\text{[149]}\)。

\textbf{地基望远镜研究}
\begin{figure}[ht]
\centering
\includegraphics[width=8cm]{./figures/9ce1392db0ded268.png}
\caption{伽利略在《星际信使》（Sidereus Nuncius）中绘制的木星及其“美第奇星”（Medicean Stars）图像} \label{fig_Jupite_10}
\end{figure}
1610 年，意大利博学者伽利略·伽利莱（Galileo Galilei）使用望远镜发现了木星的四颗最大卫星（现称伽利略卫星）。这被认为是首次利用望远镜观测地球以外的卫星。在伽利略之后仅一天，西蒙·马留斯（Simon Marius）也独立发现了木星的卫星，尽管他直到 1614 年才在书中发表其发现\(^\text{[150]}\)。然而沿用至今的是马留斯为四颗主要卫星所命名的名称：Io、Europa、Ganymede 和 Callisto。此项发现为尼古拉·哥白尼的日心说行星运动理论提供了强有力的支持；伽利略对日心说的公开支持导致他被宗教裁判所审判并定罪\(^\text{[151]}\)。

1639 年秋，来自那不勒斯的光学技师弗朗切斯科·丰塔纳（Francesco Fontana）测试了他自制的一台 22 棕榈长的望远镜，并发现了木星大气的特征性云带\(^\text{[152]}\)。

1660 年代，乔瓦尼·卡西尼（Giovanni Cassini）使用新型望远镜发现了木星大气中的斑点，观察到木星呈扁球形，并估算了其自转周期\(^\text{[153]}\)。1692 年，卡西尼注意到木星大气存在差异自转现象\(^\text{[154]}\)。

大红斑可能早在 1664 年由罗伯特·胡克（Robert Hooke）观测到，1665 年由卡西尼观测到，但这存在争议。药剂师海因里希·施瓦贝（Heinrich Schwabe）于 1831 年绘制了已知最早展示大红斑细节的图像\(^\text{[155]}\)。据记录，大红斑在 1665 年至 1708 年间数次从视野中消失，直到 1878 年才再次变得十分显眼\(^\text{[156]}\)。它在 1883 年以及 20 世纪初又被记录为再次变暗\(^\text{[157]}\)。

乔瓦尼·博雷利（Giovanni Borelli）和卡西尼均为木星卫星的运动制作了精确的表格，使人们能够预测卫星何时会从木星前方或后方经过。到 1670 年代，卡西尼注意到当木星位于太阳与地球的对侧时，这些现象发生的时间会比预计晚约 17 分钟。奥勒·罗默（Ole Rømer）据此推断光并非瞬时传播（这一结论先前被卡西尼否认）\(^\text{[51]}\)，并利用这一时间差估算了光速\(^\text{[158][159]}\)。

1892 年，E. E. 巴纳德（E. E. Barnard）利用加州利克天文台 36 英寸（910 毫米）折射望远镜观测到木星的第五颗卫星。此卫星后来被命名为 Amalthea\(^\text{[160]}\)。这是最后一颗通过望远镜视觉观测直接发现的行星卫星\(^\text{[161]}\)。在 1979 年旅行者 1 号飞掠之前，又有另外八颗木星卫星被发现\(^\text{[h]}\)。

1932 年，鲁珀特·维尔特（Rupert Wildt）在木星光谱中识别出氨与甲烷的吸收带\(^\text{[162]}\)。1938 年，木星大气中发现了三个长期存在的反气旋特征，被称为“白椭圆”。几十年来，这些椭圆在彼此靠近时并未发生合并。最终，其中两个椭圆于 1998 年合并，随后在 2000 年吸收了第三个，形成了 Oval BA\(^\text{[163]}\)。

\textbf{射电望远镜研究}

1955 年，伯纳德·伯克（Bernard Burke）和肯尼斯·富兰克林（Kenneth Franklin）发现木星会以 22.2 MHz 的频率发射无线电脉冲\(^\text{[64]:36^\text{}]}\)。这些脉冲的周期与木星的自转周期一致，他们利用这一信息更精确地测定了木星的自转速率。研究发现，来自木星的无线电脉冲有两种形式：持续数秒的长脉冲（L-bursts）以及持续不足百分之一秒的短脉冲（S-bursts）\(^\text{[164]}\)。

科学家已发现来自木星的三类无线电信号：
\begin{itemize}
\item 十米波无线电脉冲（波长达数十米）随木星自转变化，并受到木卫一与木星磁场相互作用的影响\(^\text{[165]}\)。
\item 分米波无线电辐射（波长以厘米计）首次由弗兰克·德雷克（Frank Drake）和海因·赫瓦图姆（Hein Hvatum）于 1959 年观测到\(^\text{[64]:36^\text{}]}\)。这种信号来自木星赤道周围的一个环状带，带中的电子在木星磁场的加速下产生回旋辐射\(^\text{[166]}\)。
\item 热辐射由木星大气中的热量产生\(^\text{[64]:43^\text{}]}\)。
\end{itemize}
\textbf{探测}
自 1973 年以来，自动化航天器开始探访木星，当时先驱者 10 号（Pioneer 10）探测器首次飞掠木星并传回了关于其性质和现象的突破性数据\(^\text{[167][168]}\)。前往木星的任务需要付出能量代价，其大小通常由航天器所需的速度增量（delta-v）衡量。从近地轨道进入通往木星的霍曼转移轨道需要 6.3 km/s 的速度增量\(^\text{[169]}\)，这一数值与到达近地轨道所需的 9.7 km/s 相当\(^\text{[170]}\)。通过行星飞掠提供的引力助推可以降低抵达木星所需的能量\(^\text{[171]}\)。

\textbf{飞掠任务}
\begin{figure}[ht]
\centering
\includegraphics[width=8cm]{./figures/ef4b56f8efb8c77d.png}
\caption{} \label{fig_Jupite_11}
\end{figure}
自 1973 年起，多艘航天器进行了行星飞掠机动，使其进入木星的可观测范围。先驱者号任务首次获得了木星大气及其多颗卫星的近距离图像。它们发现木星附近的辐射场远比预期强，但两艘探测器都成功在该恶劣环境中生存下来。这些航天器的轨道被用来进一步精确木星系统的质量估计。通过行星造成的无线电掩星现象，使得对木星直径及两极扁率的测量更加精确\(^\text{[59]:47^\text{}[173]}\)。

六年后，旅行者号任务极大提升了对伽利略卫星的认识，并发现了木星的行星环。它们也确认大红斑是一场反气旋风暴。图像对比显示，自先驱者号任务以来，大红斑的色调已经发生变化，从橙色转为深棕色。沿着木卫一轨道路径发现了一个由电离原子构成的环，其来源被确认是木卫一表面爆发的火山活动。当航天器从木星背面飞过时，它观测到夜侧大气中的闪电\(^\text{[59]:87^\text{}[174]}\)。

下一次遭遇木星的任务是尤利西斯号（Ulysses）太阳探测器。1992 年 2 月，它进行了一次飞掠机动以进入太阳的极轨道。在此次飞掠中，探测器研究了木星的磁层，但由于未配备摄像设备，因此无法拍摄行星图像。六年后，该探测器再次经过木星附近，但距离远得多\(^\text{[172]}\)。

2000 年，卡西尼号（Cassini）探测器在前往土星途中飞掠木星，并提供了分辨率更高的图像\(^\text{[175]}\)。

2007 年，新视野号（New Horizons）探测器在前往冥王星途中飞掠木星以获得引力助推\(^\text{[176]}\)。探测器的摄像设备测量了来自木卫一火山的等离子体输出，并详细研究了四颗伽利略卫星\(^\text{[177]}\)。

\textbf{伽利略号任务}
\begin{figure}[ht]
\centering
\includegraphics[width=8cm]{./figures/5deadf09cda53153.png}
\caption{1989 年，伽利略号在与火箭对接前的准备阶段} \label{fig_Jupite_12}
\end{figure}
第一艘进入木星轨道的航天器是伽利略号任务，它于 1995 年 12 月 7 日抵达该行星\(^\text{[71]}\)。伽利略号在轨运行超过七年，多次飞掠四颗伽利略卫星及 Amalthea。该航天器还观测到了 1994 年彗星舒梅克–列维 9 号（Shoemaker–Levy 9）与木星相撞的事件。由于高增益天线故障，任务的一些目标未能实现\(^\text{[178]}\)。

一枚 340 公斤的钛制大气探测器于 1995 年 7 月从伽利略号上释放，并于 12 月 7 日进入木星大气\(^\text{[71]}\)。它以约 2,575 公里/小时（1,600 英里/小时）的速度下降，穿越 150 公里（93 英里）的大气层\(^\text{[71]}\)，并在航天器被摧毁前共收集了 57.6 分钟的数据\(^\text{[179]}\)。伽利略轨道器本身于 2003 年 9 月 21 日被故意引导进入木星，经历了更迅速的毁灭。NASA 之所以摧毁伽利略号，是为了避免其可能坠入并污染可能存在生命的木卫二（Europa）\(^\text{[178]}\)。

此次任务的数据表明，氢在木星大气中的含量高达 90\%\(^\text{[71]}\)。记录的温度超过 300 °C（572 °F），在探测器被汽化前测得的风速超过 644 公里/小时（>400 英里/小时）\(^\text{[71]}\)。

\textbf{朱诺号任务}
\begin{figure}[ht]
\centering
\includegraphics[width=8cm]{./figures/3db8b370dca0782e.png}
\caption{2011 年，朱诺号在旋转测试台上准备接受测试} \label{fig_Jupite_13}
\end{figure}
NASA 的朱诺号任务于 2016 年 7 月 4 日抵达木星，目标是在极轨道上对木星进行详细研究。该航天器原本计划在二十个月内绕木星运行三十七圈\(^\text{[78][180][181]}\)。在任务期间，航天器将暴露于来自木星磁层的高强度辐射，这可能导致部分仪器失效\(^\text{[182]}\)。2016 年 8 月 27 日，朱诺号完成了首次飞掠木星，并传回了历史上首批木星北极图像\(^\text{[183]}\)。

朱诺号在 2018 年 7 月预算任务结束前共完成了 12 次轨道运行\(^\text{[184]}\)。当年 6 月，NASA 将任务计划延长至 2021 年 7 月；同年 1 月，任务再次延长至 2025 年 9 月，包含四次卫星飞掠：一次木卫三（Ganymede）、一次木卫二（Europa）以及两次木卫一（Io）\(^\text{[185][186]}\)。当朱诺任务结束时，将进行受控脱轨并在木星大气中解体，以避免撞击并污染木星的卫星\(^\text{[187]}\)。

\textbf{取消的任务与未来计划}

科研界对研究木星较大的冰质卫星非常感兴趣，这些卫星可能拥有地下液态海洋\(^\text{[188]}\)。然而经费困难延缓了任务进展，NASA 的木星冰卫星轨道器（JIMO）计划于 2005 年被取消\(^\text{[189]}\)。随后提出了一项 NASA/ESA 联合任务 EJSM/Laplace，计划于约 2020 年发射。EJSM/Laplace 将由 NASA 主导的木卫二轨道器（Jupiter Europa Orbiter）和 ESA 主导的木卫三轨道器（Jupiter Ganymede Orbiter）组成\(^\text{[190]}\)。但 ESA 于 2011 年 4 月正式终止与 NASA 的合作，原因是 NASA 的预算问题及其对任务时间表的影响。取而代之的是，ESA 决定推进一项完全由欧洲自主执行的任务，以参加其 L1 Cosmic Vision 选拔\(^\text{[191]}\)。这些计划后来实现为欧洲航天局的“木星冰卫星探测器”（JUICE），已于 2023 年 4 月 14 日发射\(^\text{[192]}\)，随后是 NASA 的木卫二探测器 Europa Clipper，于 2024 年 10 月 14 日发射\(^\text{[193]}\)。

其他提出的任务包括中国国家航天局的“天问四号”（Tianwen-4），计划于 2035 年前后向木星系统甚至可能是木卫四（Callisto）发射轨道器\(^\text{[194]}\)，以及 CNSA 的“穿越星际快递”（Interstellar Express）\(^\text{[195]}\)和 NASA 的 Interstellar Probe\(^\text{[196]}\)，两者都将利用木星引力助推抵达日球层边缘。
\subsection{卫星}
木星目前已知拥有 97 颗天然卫星\(^\text{[8]}\)，随着望远镜观测的不断增加，这一数字很可能继续上升\(^\text{[197]}\)。在这些卫星中，仅有 16 颗直径超过 10 公里\(^\text{[198]}\)。其中最大的四颗卫星——伽利略卫星（Galilean moons）——依大小顺序分别为木卫三（Ganymede）、木卫四（Callisto）、木卫一（Io）和木卫二（Europa），在晴朗夜晚使用双筒望远镜即可从地球看到它们\(^\text{[199]}\)。
\subsubsection{伽利略卫星}
由伽利略发现的四颗卫星——Io、Europa、Ganymede 和 Callisto——是太阳系中最大的一批卫星。Io、Europa 与 Ganymede 的轨道构成一种称为拉普拉斯共振（Laplace resonance）的模式：在木卫一完成四次环绕木星的轨道周期时，木卫二正好完成两次，木卫三则正好完成一次。该共振导致三颗大型卫星之间的引力作用在它们每次绕木星的轨道中的同一点对其施加额外拉力，从而将其轨道拉长为椭圆形。另一方面，来自木星的潮汐力则倾向于使它们的轨道趋于圆形\(^\text{[200]}\)。

轨道偏心率使三颗卫星的形状定期发生弯曲：接近木星时，它们被引力拉伸，而远离木星时又恢复为更接近球形。潮汐弯曲产生的摩擦在卫星内部生成热量\(^\text{[201]}\)。这在木卫一上表现得最为显著，因为它受到最强烈的潮汐力\(^\text{[201]}\)；同时，木卫二表面地质的年轻性（这表明其表面曾被近期重新塑造）也在较低程度上体现了这一效应\(^\text{[202]}\)。
\begin{figure}[ht]
\centering
\includegraphics[width=10cm]{./figures/3cf8c93e3bbfd3fa.png}
\caption{} \label{fig_Jupite_14}
\end{figure}
\begin{figure}[ht]
\centering
\includegraphics[width=10cm]{./figures/21f593e7fc447743.png}
\caption{伽利略卫星Io,欧罗巴,木卫三,和木卫四(以增加与木星距离的顺序) 在假色} \label{fig_Jupite_15}
\end{figure}
\subsubsection{分类}
木星的卫星最初被根据其相似的轨道要素分为四组，每组四颗\(^\text{[203]}\)。自 1999 年以来大量小型外层卫星的发现使这一结构变得更加复杂。木星的卫星目前被划分为多个不同的群组，但已知有两颗卫星（Themisto 与 Valetudo）不属于任何群组\(^\text{[204]}\)。

最内侧的八颗规则卫星具有近乎圆形的轨道，且轨道平面接近木星赤道，被认为与木星共同形成。而其余卫星均为不规则卫星，被认为是被捕获的小行星或被捕获小行星的碎片。这些不规则卫星在各自群组内部可能具有共同起源，或许是一颗更大的卫星或被捕获的天体发生破裂后形成的\(^\text{[205][206]}\)。
\begin{figure}[ht]
\centering
\includegraphics[width=14.25cm]{./figures/2f0f35329bb1b88c.png}
\caption{} \label{fig_Jupite_16}
\end{figure}
\subsection{与太阳系的相互作用}

作为八大行星中质量最大的行星，木星的引力影响深刻塑造了太阳系。除水星之外，其余所有行星的轨道面都更接近木星的轨道平面，而不是太阳的赤道平面。小行星带中的科克伍德空隙主要由木星造成\(^\text{[211]}\)，并且木星可能也在太阳系内侧历史上的“晚期重轰炸”（Late Heavy Bombardment）事件中发挥了重要作用\(^\text{[212]}\)。

除了拥有众多卫星之外，木星的引力场还控制着大量小行星，这些小行星聚集在木星绕太阳轨道上领先与落后的拉格朗日点周围。这些天体被称为木星特洛伊小行星（Trojan asteroids），并以《伊利亚特》为名分为“希腊营”与“特洛伊营”。其中首颗被发现的是 588 Achilles，由马克斯·沃尔夫（Max Wolf）于 1906 年发现；此后又发现了两千多颗\(^\text{[213]}\)。最大的特洛伊小行星是 624 Hektor\(^\text{[214]}\)。

所谓“木星族彗星”（Jupiter family）指半长轴小于木星轨道半长轴的彗星；大多数短周期彗星都属于这一族群。木星族彗星被认为起源于海王星轨道外侧的柯伊伯带。在与木星的近距离遭遇中，它们的轨道会受到扰动而进入更短周期的轨道，并通过太阳与木星的周期性引力作用使轨道逐渐变为圆形\(^\text{[215]}\)。
\subsubsection{撞击}
\begin{figure}[ht]
\centering
\includegraphics[width=8cm]{./figures/dc70a8a94d1bf9b1.png}
\caption{棕色斑点标示出彗星舒梅克–列维 9 号（Shoemaker–Levy 9）在木星上的撞击地点} \label{fig_Jupite_17}
\end{figure}
由于其巨大的引力势阱以及位于太阳系内侧邻近的位置，木星被称为“太阳系的吸尘器”\(^\text{[216]}\)。木星遭受的撞击——例如彗星撞击——比太阳系中任何其他行星都多\(^\text{[217]}\)。例如，木星受到的小行星和彗星撞击次数约为地球的 200 倍\(^\text{[71]}\)。科学家过去认为木星在一定程度上保护了内太阳系免受彗星轰击\(^\text{[71]}\)。然而，2008 年的电脑模拟表明，木星并不会使穿越内太阳系的彗星总体数量减少，因为它通过引力扰动把彗星轨道向内拉的频率，与其吞噬或抛出彗星的频率大致相当\(^\text{[218]}\)。这一议题至今仍在科学界存在争议：一些人认为木星将彗星从柯伊伯带引向地球，而另一些人则认为木星保护地球免受来自奥尔特云的彗星影响\(^\text{[219]}\)。

1994 年 7 月，彗星舒梅克–列维 9 号（Shoemaker–Levy 9）与木星相撞\(^\text{[220][221]}\)。此次撞击被全球多个天文台密切观测，包括哈勃空间望远镜与伽利略号探测器\(^\text{[222][223][224]}\)。这一事件也受到媒体的广泛报道\(^\text{[225]}\)。

对早期天文记录和绘图的调查揭示了 1664 年至 1839 年间可能出现的八个撞击观测案例。然而，1997 年的一项回顾研究认定这些观测几乎没有可能是真实撞击的结果。该团队的进一步调查揭示，天文学家乔瓦尼·卡西尼（Giovanni Cassini）在 1690 年发现的一处暗色表面特征可能是一次撞击留下的痕迹\(^\text{[226]}\)。
\subsection{文化中的木星}
\begin{figure}[ht]
\centering
\includegraphics[width=8cm]{./figures/c61eb8602d5a23f5.png}
\caption{1550 年版圭多·博纳蒂（Guido Bonatti）《天文学书》（Liber Astronomiae）中的木星木刻插图} \label{fig_Jupite_18}
\end{figure}
\begin{figure}[ht]
\centering
\includegraphics[width=8cm]{./figures/ce5eed22ea6de744.png}
\caption{云中持雷霆与权杖的朱庇特—宙斯（Jupiter–Zeus），伴随其鹰的形象。赫库兰尼姆（Herculaneum）壁画，公元 1–37 年} \label{fig_Jupite_19}
\end{figure}
自古以来，人类就已知晓木星的存在。它在夜空中肉眼可见，并且在太阳较低时也可以在白天看到\(^\text{[227]}\)。对于巴比伦人而言，这颗行星代表着他们的神祇马尔杜克（Marduk）\(^\text{[228]}\)，自汉谟拉比时期起即为其万神殿的主神\(^\text{[229]}\)。他们利用木星沿黄道大约 12 年的轨道周期来确定其黄道十二宫的星座体系\(^\text{[228]}\)。

木星在希腊神话中的名称是宙斯（Ζεύς），亦称为迪亚斯（Δίας），而现代希腊语仍沿用该行星名\(^\text{[230]}\)。古希腊人称其为“Phaethon”（Φαέθων），意为“闪耀者”或“燃烧之星”\(^\text{[231][232]}\)。荷马时期的宙斯神话与近东诸神之间存在明显相似性，包括闪米特神祇 El 与 Baal、苏美尔神 Enlil，以及巴比伦神 Marduk\(^\text{[233]}\)。木星与希腊神宙斯的对应关系来自近东文化的影响，在公元前四世纪已完全确立，如柏拉图及其同时代人的《Epinomis》中所记载\(^\text{[234]}\)。

朱庇特（Jupiter）是宙斯的罗马对应神祇，也是罗马神话中的主神。罗马人最初称木星为“Iuppiter Stella”（朱庇特之星），因为他们认为该行星神圣并属于朱庇特。这一名称来源于原始印欧语的呼格复合词 *Dyēu-pəter（主格：*Dyēus-pətēr，意为“天空之父神”或“白昼之父神”）\(^\text{[235]}\)。作为罗马万神殿的至高神，朱庇特是雷、电与风暴之神，同时被称为光明与天空之神\(^\text{[236]}\)。

在吠陀占星学中，印度占星家以 Brihaspati（众神的宗教导师）为木星命名，并称之为“Guru”，意为“教师”\(^\text{[237][238]}\)。在中亚突厥神话中，木星被称为 Erendiz 或 Erentüz，来自 eren（含义不确定）与 yultuz（“星”）两词。突厥人认为木星的轨道周期为 11 年 300 天，并相信某些社会与自然现象与 Erentüz 在天空中的运行有关\(^\text{[239]}\)。在中国、越南、韩国与日本，木星被称为“木星”，基于中国五行学说（Chinese: 木星; pinyin: mùxīng）\(^\text{[240][241][242]}\)。在中国，它亦被称为“岁星”（Chinese: 歲星; pinyin: suìxīng），因中国天文学家注意到其每年（经过校正）在黄道中“跳”至下一个星座。在某些古代中国文献中，年份在原理上是按照木星十二次星座周期（岁星纪年）命名的\(^\text{[143]}\)。
\subsection{另见}
\begin{itemize}
\item Outline of Jupiter – 太阳系第五颗行星木星概览  
\item Eccentric Jupiter – 绕其恒星运行轨道偏心的类木行星  
\item Hot Jupiter – 近距离环绕恒星运行的高质量行星  
\item Super-Jupiter – 质量大于木星的一类行星  
\item Jovian–Plutonian gravitational effect – “木星–冥王星引力效应”（一则天文骗局）  
\item List of gravitationally rounded objects of the Solar System – 太阳系中因引力而呈球形的天体列表  
\end{itemize}
\subsection{注释}\\
a.光学滤镜的波长为 658 nm、502 nm 和 395 nm\(^\text{[1]}\)，分别大致对应红、绿、蓝三色。\\  
b.木星的大气与外观不断变化，因此今日其外观可能与拍摄该图像时不尽相同。然而此图像中仍显示一些持久特征，例如著名的大红斑（图像右下位置极为显著）以及木星典型的带状结构。\\  
c.指 1 bar 大气压所对应的大气层高度。\\  
d.基于 1 bar 大气压层以内的体积估算。\\  
e.100 ×（赤道半径 − 极半径）/（赤道半径）＝ 100 ×（71492 − 66854）/ 71492 ＝ 6.487\%。\\  
f.参见例如："IAUC 2844: Jupiter; 1975h". International Astronomical Union. October 1, 1975. Retrieved October 24, 2010. 该词至少自 1966 年起使用，参见："Query Results from the Astronomy Database". Smithsonian/NASA. Retrieved July 29, 2007. \\
g.大致与糖浆（syrup USP）的密度相当\(^\text{[43]}\)。 \\ 
h.详情与引用见 Moons of Jupiter 条目。
\subsection{参考资料}
\begin{enumerate}
\item “木星蛋白石2024”。NASA.2024。已检索5月10日，2025。
\item 辛普森，J。A.;韦纳，E。S。C.(1989)。“木星”。牛津英语词典。第8卷 (第二版)。克拉伦登.ISBN 978-0-19-861220-9。
\item 威廉姆斯，大卫R.(2021年12月23日)。“木星概况介绍”。NASA。已存档从原始的2019年12月29日。已检索10月13日，2017。
\item 塞利格曼，考特尼。“轮换周期和日长”。已存档从原来的2018年9月29日。已检索8月13日，2009年。
\item 西蒙，J.L.;布列塔尼翁，P.；查普朗；查普朗-图兹，M.；弗朗库，G.；拉斯卡，J。(1994年2月)。“月球和行星的进动公式和平均元素的数值表达式”。天文学和天体物理学。282(2):663-683。Bibcode:1994A & A。282 .. 663S。
\item 苏阿米，D.；Souchay, J.(2012年7月)。“太阳系的不变平面”。天文学与天体物理学。543: 11。Bibcode:2012A & A...543A.133S。doi:10.1051/0004-6361/201219011。A133.
\item “HORIZONS Planet-center批量呼叫2023年1月的近日点”。ssd.jpl.nasa.gov(木星的行星中心 (599) 的近日点发生在2023年1月21日，在rdot从负到正的翻转期间，在4.9510113au)。NASA/JPL。已存档从原始的2021年9月7日。已检索9月7日，2021。
\item “行星卫星发现情况”。JPL太阳能系统动力学。NASA.2025年4月30日已存档从原始的2021年9月27日。已检索4月30日，2025。
\item Seidelmann，P。肯尼斯; 阿基纳，布伦特A；迈克尔·F·阿赫恩；阿尔伯特·R·康拉德；Consolmagno, Guy J.;丹尼尔·赫斯特罗弗; 詹姆斯·L·希尔顿；乔治·克拉辛斯基；格雷戈里·诺伊曼；尤尔根·奥伯斯特; 菲利普·J·斯托克；爱德华·F·特德斯科；大卫·J·索伦；彼得·C·托马斯；威廉姆斯，伊万·P。(2007)。“IAU/IAG制图坐标和旋转要素工作组的报告: 2006年”。天体力学与动力天文学。98(3):155-180。Bibcode:2007CeMDA .. 98 .. 155S。doi:10.1007/s10569-007-9072-y。ISSN0923-2958。 
\item 德佩特，Imke; 利绍尔，杰克J。(2015)。行星科学(第二次更新版本)。纽约: 剑桥大学出版社。第250页。ISBN 978-0-521-85371-2。已检索8月17日，2016。
\item "天体动力学常数"。JPL太阳能系统动力学。2009年2月27日已存档从原始的2019年3月21日。已检索8月8日，2007。
\item “NASA: 太阳系探索: 行星: 木星: 事实与数据”。solarsystem.nasa.gov。2011年6月2日。存档自原来的2011年9月5日。已检索10月15日，2024。
\item Ni, D.(2018)。“朱诺数据对木星内部的经验模型”。天文学与天体物理学。613: a32。Bibcode:2018A & A...613A..32N。doi:10.1051/0004-6361/201732183。
\item 阿基纳，B。A.;阿克顿，C。H、; A'Hearn，M。F.;康拉德，A.；Consolmagno, G.J.;达克斯伯里，T.；赫斯特罗弗，D.；希尔顿，J。L.;柯克，R.L.;Klioner, S.A.;麦卡锡，D.；米奇，K.；奥伯斯特，J.；平，J.；Seidelmann，P。K。(2018)。“IAU制图坐标和旋转要素工作组的报告: 2015年”。天体力学与动力天文学。130(3): 22。Bibcode:2018 CeMDA.130. .. 22A。doi:10.1007/s10569-017-9805-5。ISSN0923-2958。 
\item 李，李明; 江，X；西，R。A.;Gierasch，P。J.;Perez-Hoyos, S.;桑切斯-拉维加，A.；弗莱彻，L。N.;福特尼，J.J.;诺尔斯，B.；波科，C。C.;贝恩斯，K。H、; 弗莱，P。M、；Mallama, A.;阿克特伯格，R。K.;西蒙，A.A.;尼克松，C。A.;奥顿，G.美国；Dyudina，美国。A.;Ewald, S.P.;Schmude, R.W。(2018)。“更少吸收的太阳能和更多的木星内部热量”。自然通讯。9(1): 3709。Bibcode:2018NatCo...9.3709L。doi:10.1038/s41467-018-06107-2。PMC6137063。PMID30213944。  
\item 马拉玛，安东尼; 克鲁布塞克，布鲁斯; 巴夫洛夫，赫里斯托 (2017)。“行星的综合宽带震级和反照率，适用于系外行星和第九行星”。伊卡洛斯。282:19-33。arXiv:1609.05048。Bibcode:2017Icar .. 282...19米。doi:10.1016/j.icarus.2016.09.023。S2CID119307693。 
\item Mallama, A.;希尔顿，J。L.(2018)。“计算天文历书的视在行星星等”。天文学与计算。25:10-24.arXiv:1808.01973。Bibcode:2018A & C .... 25...10米。doi:10.1016/j.ascom.2018.08.002。S2CID69912809。 
\item “百科全书-最聪明的身体”。IMCCE。已存档从原来的2023年7月24日。已检索5月29日，2023。
\item Bjoraker, G.L.;王，M.H.; 德佩特，I.；Ádámkovics, M.(2015年9月)。“使用Keck观测光谱分辨的线形揭示了木星的深层云结构”。天体物理学杂志。810(2): 10。arXiv:1508.04795。Bibcode:2015年apj。810。122B。doi:10.1088/0004-637X/810/2/122。S2CID55592285。122. 
\item 亚历山大，雷切尔 (2015)。太阳系天体的神话、符号和传说。帕特里克·摩尔实用天文学系列第177卷。纽约: 施普林格。pp. 141-159。Bibcode:2015年msl。书 ..... A。doi:10.1007/978-1-4614-7067-0。ISBN 978-1-4614-7066-3。
\item “天文物体的命名”。国际天文学联合会。存档自原来的2013年10月31日。已检索3月23日，2022年。
\item 亚历山大·琼斯 (1999)。来自Oxyrhynchus的天文纸莎草。美国哲学学会.pp. 62-63.ISBN 978-0-87169-233-7。现在可以将五个行星中至少四个的中世纪符号追溯到一些最新的纸莎草纸星座中出现的形式 ([P、氧。] 4272、4274、4275 [...])。对于木星来说，这是一个明显的字母组合，源自希腊名字的首字母。
\item 马德，A。S。D.(1934年8月)。“行星符号的起源”。天文台。57:238-247。Bibcode:1934Obs。57。238米。
\item 哈珀，道格拉斯。“朱庇特”。在线词源词典。已存档从原来的2022年3月23日。已检索3月22日，2022年。
\item "快乐"。Dictionary.com。已存档从原来的2012年2月16日。已检索7月29日，2007。
\item 托马斯·S·克鲁伊耶；伯克哈特，克里斯托夫; 布德，杰里特; 克莱恩，托斯滕 (2017年6月)。“从陨石的独特遗传学和形成时间推断出木星的年龄”。美国国家科学院院刊。114(26):6712-6716。Bibcode:2017PNAS..114.6712K。doi:10.1073/pnas.1704461114。PMC5495263。PMID28607079。  
\item 博斯曼，A.D.;克里德兰，A.J.;米格尔，Y.(2019年12月)。“木星形成为N2冰线周围的卵石堆”。天文学与天体物理学。632: 5。arXiv:1911.11154。Bibcode:2019A & A...632L..11B。doi:10.1051/0004-6361/201936827。S2CID208291392。L11。 
\item D'Angelo, G.;魏登斯林，S。J.;利绍尔，J。J.;博登海默，P。(2021)。“木星的生长: 气体和固体盘中的形成以及到现在的演化”。伊卡洛斯。355114087。arXiv:2009.05575。Bibcode:2021年Icar .. 35514087D。doi:10.1016/j.icarus.2020.114087。S2CID221654962。 
\item 沃尔什，K.J.;莫比德利，A.；Raymond, S.N.;O'Brien, D.P.;曼德尔，A.M。(2011)。“木星早期气体驱动迁移对火星的低质量”。性质。475(7355):206-209。arXiv:1201.5177。Bibcode:2011自然475..206w。doi:10.1038/自然10201。PMID21642961。S2CID4431823。  
\item 康斯坦丁·巴蒂金(2015)。木星在太阳系内部早期演化中的决定性作用。美国国家科学院院刊。112(14):4214-4217。arXiv:1503.06945。Bibcode:2015PNAS..112.4214B。doi:10.1073/pnas.1423252112。PMC4394287。PMID25831540。  
\item Chametla，ra ú l O; D'Angelo，Gennaro; 雷耶斯-鲁伊斯，毛里西奥; 桑切斯-萨尔塞多，F哈维尔 (2020年3月)。“在气态原行星盘的平均运动共振中捕获和迁移木星和土星”。皇家天文学会月刊。492(4):6007-6018。arXiv:2001.09235。doi:10.1093/mnras/staa260。
\item 小海施，K。E.;拉达，E。A.;拉达，C。J.(2001)。“年轻星团中的圆盘频率和寿命”。天体物理学杂志。553(2):153-156。arXiv:astro-ph/0104347。Bibcode:2001ApJ...553L.153H。doi:10.1086/320685。S2CID16480998。 
\item Fazekas，安德鲁 (2015年3月24日)。“观察: 木星，破坏早期太阳系的球”。国家地理。存档自原来的2017年3月14日。已检索4月18日，2021。
\item 北卡罗来纳州祖贝；尼莫，F.；费舍尔，R.；雅各布森，S.(2019年)。“Hf-w同位素演化对大粘性场景中地球行星形成时间尺度和平衡过程的限制”。地球和行星科学快报。522(1):210-218.arXiv:1910.00645。Bibcode:2019E & PSL.522..210Z。doi:10.1016/j.epsl.2019.07.001。PMC7339907。PMID32636530。S2CID199100280。   
\item D'Angelo, G.;Marzari, F.(2012)。“木星和土星在演化的气态盘中向外迁移”。天体物理学杂志。757(1): 50 (23页)。arXiv:1207.2737。Bibcode:2012年apj... 757...50D。doi:10.1088/0004-637X/757/1/50。S2CID118587166。 
\item 皮拉尼，S.；约翰森，A.；比特施，B.；穆斯蒂尔，A.J.；Turrini, D.(2019年3月)。“行星迁移对早期太阳系小天体的影响”。天文学与天体物理学。623: a169。arXiv:1902.04591。Bibcode:2019A & A...623A.169P。doi:10.1051/0004-6361/201833713。
\item “木星的未知旅程揭晓”。科学日报。隆德大学。2019年3月22日。已存档从原始的2019年3月22日。已检索3月25日，2019。
\item 哈罗德·F·利维森；莫比德利，亚历山德罗; 范·拉尔霍芬，克里斯塔; 戈麦斯，R。(2008)。“天王星和海王星轨道动态不稳定期间柯伊伯带结构的起源”。伊卡洛斯。196(1):258-273.arXiv:0712.0553。Bibcode:2008Icar .. 196 .. 258L。doi:10.1016/j.icarus.2007.11.035。S2CID7035885。 
\item 康斯坦丁·巴蒂金; 迈克尔·E·布朗。(2016)。“太阳系中一颗遥远的巨行星的证据”。天文杂志。151(2): 22。arXiv:1601.05438。Bibcode:2016AJ....151...22B。doi:10.3847/0004-6256/151/2/22。
\item Öberg, K.I.;华兹华斯，R.(2019年)。木星的组成表明它的核心组装在n_ {2} 雪线的外部。。天文杂志。158(5)。arXiv:1909.11246。doi:10.3847/1538-3881/ab46a8。S2CID202749962。 
\item Öberg, K.I.;华兹华斯，R.(2020)。勘误表: “木星的组成表明其核心组装在N2雪线的外部”"。天文杂志。159(2): 78。doi:10.3847/1538-3881/ab6172。S2CID 214576608。
\item 德内克，爱德华·J。(2020年1月7日)。摄政考试和答案: 地球科学-物理环境2020。巴伦教育系列。第419页。ISBN 978-1-5062-5399-2。
\item 詹姆斯·斯沃布里克 (2013)。制药技术百科全书。第6卷。CRC出版社。第3601页。ISBN 978-1-4398-0823-8。已存档从原来的2023年3月26日。已检索3月19日，2023。糖浆USP (1.31克/厘米3)
\item 克拉本·沃尔特·艾伦；考克斯，亚瑟·N。(2000)。艾伦的天体物理量。斯普林格.pp. 295-296。ISBN 978-0-387-98746-0。已存档从原来的2023年2月21日。已检索3月18日，2022年。
\item Polyanin, Andrei D.;Chernoutsan，阿列克谢 (2010年10月18日)。数学，物理和工程科学的简明手册。CRC出版社。第1041页。ISBN 978-1-4398-0640-1。
\item 吉洛，特里斯坦; 戈蒂埃，丹尼尔; 哈伯德，威廉B。(1997年12月)。“注: 伽利略测量和内部模型对木星组成的新限制”。伊卡洛斯。130(2):534-539。arXiv:astro-ph/9707210。Bibcode:1997年Icar。。130。534克。doi:10.1006/icar.1997.5812。S2CID5466469。 
\item 弗兰·巴格纳尔; 蒂莫西·E·道林；麦金农，威廉·B，编辑。(2006)。木星: 行星、卫星和磁层。剑桥大学出版社。pp. 59-75。ISBN 0-521-03545-7。
\item 金，S.J.;考德威尔，J.；Rivolo, A.R、；瓦格纳，R。(1985)。“木星III上的红外极地增亮。来自旅行者1号虹膜实验的光谱法 “。伊卡洛斯。64(2):233-248。Bibcode:1985年伊卡。64。233K。doi:10.1016/0019-1035(85)90201-5。
\item Vdovichenko，V。D.;卡里莫夫，A。M、；基里延科，G。A.;Lysenko, P.G.;Tejfel'，V。G.;菲利波夫，V.A.;Kharitonova, G.A.;Khozhenets, A.P。(2021)。“木星上弱分子吸收带行为的带状特征”。太阳系研究。55(1):35-46.Bibcode:2021年索斯里。。55。35伏。doi:10.1134/S003809462101010X。S2CID255069821。 
\item 戈蒂埃，D.；康拉斯，B.；弗拉萨尔，M.；哈内尔，R.；昆德，V.；Chedin, A.;斯科特，N.(1981)。“来自旅行者号的木星氦丰度”。地球物理研究杂志。86(A10):8713-8720。Bibcode:1981JGR....86.8713G。doi:10.1029/JA086iA10p08713。高密度脂蛋白:2060/19810016480。S2CID122314894。 
\item 昆德，V。G.;弗拉萨尔，F。M、；詹宁斯，D。E.;Bézard, B.;斯特罗贝尔，D。F.; 等人 (2004年9月10日)。“卡西尼热红外光谱实验中的木星的大气成分”。科学。305(5690):1582-1586。Bibcode:2004Sci...305.1582K。doi:10.1126/science.1100240。PMID15319491。S2CID45296656。  
\item “太阳星云超市”(PDF)。nasa.gov。2017年12月。已存档(PDF)从原来的2023年7月17日。已检索7月10日，2023。
\item 尼曼，H. B.；Atreya, S.K.;Carignan, G.R、；多纳休，T。M、；哈伯曼，J。A.; 等人 (1996年)。“伽利略探针质谱仪: 木星的大气组成”。科学。272(5263):846-849。Bibcode:1996Sci。272。846N。doi:10.1126/science.272.5263.846。PMID8629016。S2CID3242002。  
\item 冯·扎恩，美国；亨顿，D.M、；Lehmacher，G。(1998)。“木星大气中的氦: 伽利略探测器氦干涉仪实验的结果”。地球物理研究杂志。103(E10):22815-22829。Bibcode:1998JGR...10322815V。doi:10.1029/98JE00695。
\item 史蒂文森，大卫J.(2020年5月)。“朱诺揭示的木星内部”。地球与行星科学年度评论。48:465-489。Bibcode:2020AREPS..48..465S。doi:10.1146/annurev-earth-081619-052855。S2CID212832169。 
\item 英格索尔，A.P.;哈默尔，H. B.；Spilker, T.R、；Young, R.E.(2005年6月1日)。“外行星: 冰巨人”(PDF)。月球与行星研究所。已存档(PDF)从原来的2022年10月9日。已检索2月1日，2007。
\item 马克·霍夫斯塔特 (2011)，“冰巨人，天王星和海王星的大气”(PDF),白皮书行星科学十年调查,美国国家研究委员会,pp。 1-2、已存档(PDF)从原来的2023年7月17日,检索到1月18日，2015
\item 麦克道格尔，道格拉斯W.(2012)。“离家近的双星系统: 月球和地球如何绕轨道运行”。牛顿引力。本科物理讲义。施普林格纽约。pp. 193-211.doi:10.1007/978-1-4614-5444-1_10。ISBN 978-1-4614-5443-4。重心距太阳中心743，000公里。太阳的半径为696，000公里，因此距地表47，000公里。
\item 伯吉斯，埃里克 (1982)。木星: 奥德赛到一个巨人。纽约: 哥伦比亚大学出版社。ISBN 978-0-231-05176-7。
\item 舒，弗兰克·H (1982)。物理宇宙: 天文学导论。天文学丛书 (第12版)。大学科学书籍。 426。ISBN 978-0-935702-05-7。
\item 安德鲁·M·戴维斯；Turekian, Karl K.(2005)。陨石、彗星和行星。地球化学论文.第1卷。爱思唯尔。第624页。ISBN 978-0-08-044720-9。
\item 施耐德，让 (2009)。“太阳系外行星百科全书: 交互式目录”。太阳系外行星百科全书。已存档从原来的2023年10月28日。已检索8月9日，2014。
\item 冯，法博; 巴特勒，R。保罗; 等人 (2022年8月)。“三维选择167个次恒星伴星到附近恒星”。天体物理学杂志增刊系列。262(21): 21。arXiv:2208.12720。Bibcode:2022年ApJS。。262。21F。doi:10.3847/1538-4365/ac7e57。S2CID251864022。 
\item 埃尔金斯-坦顿，琳达T.(2011)。木星和土星(修订版)。纽约: 切尔西之家。ISBN 978-0-8160-7698-7。
\item 欧文，帕特里克 (2003)。我们的太阳系的巨型行星: 大气，组成和结构。施普林格科学与商业媒体。第62页。ISBN 978-3-540-00681-7。已存档从原来的2024年6月19日。已检索4月23日，2022年。
\item 欧文，帕特里克G.J.(2009) [2003]。我们的太阳系的巨型行星: 大气，组成和结构(第二版)。斯普林格。第4页。ISBN 978-3-642-09888-8。已存档从原来的2024年6月19日。已检索3月6日，2021。据估计，木星的半径目前正在缩小约1毫米/年。
\item 吉洛，特里斯坦; 史蒂文森，大卫·J；威廉·B·哈伯德；索门，迪迪埃 (2004年)。第三章木星内部在弗兰的巴格纳尔; 道林，蒂莫西·E；威廉·B·麦金农 (主编)。木星: 行星、卫星和磁层。剑桥大学出版社。ISBN 978-0-521-81808-7。
\item 博登海默，P。(1974)。“木星的早期演化的计算”。伊卡洛斯。23.23(3):319-325。Bibcode:1974年Icar。23 .. 319B。doi:10.1016/0019-1035(74)90050-5。
\item “木星以前是现在的两倍，磁场更强。。加州理工学院。2025年5月20日。已检索10月30日，2025。
\item 西格，S.；库什内尔，M.；Hier-majumder，C。A.;Militzer, B.(2007)。“固体系外行星的质量-半径关系”。天体物理学杂志。669(2):1279-1297。arXiv:0707.2895。Bibcode:2007ApJ...669.1279S。doi:10.1086/521346。S2CID8369390。 
\item 3宇宙如何运作。卷。木星: 驱逐舰还是救世主？。探索频道。2014。
\item 特里斯坦·吉洛 (1999)。“太阳系内外巨型行星的内部”(PDF)。科学。286(5437):72-77.Bibcode:1999Sci...286...72G。doi:10.1126/science.286.5437.72。PMID10506563。已存档(PDF)从原来的2022年10月9日。已检索4月24日，2022年。  
\item 伯罗斯，亚当; 哈伯德，W。B.;Lunine, J.I.;詹姆斯·利伯特 (2001年7月)。“褐矮星和太阳系外巨行星的理论”。现代物理学评论。73(3):719-765。arXiv:astro-ph/0103383。Bibcode:2001年RvMP。73 .. 719B。doi:10.1103/RevModPhys.73.719。S2CID204927572。因此，太阳能金属度和Y的HBMMα= 50.25等于0.07 - 0.074M☉,...而零金属度下的HBMM为0.092M☉ 
\item 亚历山大·冯·博蒂彻; 阿莫里·H·M·特里奥德。J.;奎洛兹，迪迪尔; 吉尔，萨姆; 伦德尔，莫妮卡; 德尔雷兹，莱蒂西亚; 安德森，大卫·R。；科利尔·卡梅隆，安德鲁; 法埃迪，弗朗西斯卡; 吉隆，米歇尔; 戈麦斯·马奎奥·秋，伊伦; 赫布，莱斯利; 赫利尔，科埃尔; 杰辛，伊曼纽尔; 马克斯泰德，皮埃尔·F。L.;大卫·V·马丁；佩佩，弗朗切斯科; Pollacco，唐; 塞格兰桑，达米安; 斯莫利，巴里; Udry，st é phane; West，理查德 (2017年8月)。“EBLM项目。三。一颗土星大小的低质量恒星，处于氢燃烧极限。天文学与天体物理学。604: 6。arXiv:1706.08781。Bibcode:2017A & A...604L...6V。doi:10.1051/0004-6361/201731107。S2CID54610182。L6。 
\item Smoluchowski，R。(1971)。“木星和土星的金属内饰和磁场”。天体物理学杂志。166: 435。Bibcode:1971年4月。166 .. 435S。doi:10.1086/150971。
\item Wahl, S.M、；威廉·B·哈伯德；军事，B.；吉洛，特里斯坦; 米格尔，Y；Movshovitz, N.;Kaspi, Y.;赫尔，R；里斯，D.；加兰蒂，E.；莱文，S.；康纳尼，J。E.;博尔顿，S.J.(2017)。“比较木星内部结构模型与朱诺重力测量和稀释核心的作用”。地球物理研究快报。44(10):4649-4659。arXiv:1707.01997。Bibcode:2017乔治。。44.4649W。doi:10.1002/2017GL073160。
\item 刘尚飞; 等 (2019年8月15日)。“木星的稀释核心的形成受到巨大的影响”。性质。572(7769):355-357。arXiv:2007.08338。Bibcode:2019 Natur.572..355L。doi:10.1038/s41586-019-1470-2。PMID31413376。S2CID199576704。  
\item Chang，Kenneth (2016年7月5日)。“美国宇航局的朱诺号飞船进入木星的轨道”。纽约时报。已存档从2019年5月2日的原始。已检索7月5日，2016。
\item 史蒂文森，D。J.;博登海默，P.；利绍尔，J。J.;D'Angelo, G.(2022年)。“木星形成和演化过程中可凝结成分与h-he的混合”。行星科学杂志。3(4): id.74。arXiv:2202.09476。Bibcode:2022PSJ.....3...74S。doi:10.3847/PSJ/ac5c44。S2CID247011195。 
\item 内疚，T。(2019年)。“有迹象表明木星受到巨大撞击的影响”。性质。572(7769):315-317。Bibcode:2019 natur.572.315克。doi:10.1038/d41586-019-02401-1。PMID31413374。 
\item Trachenko, K.;Brazhkin, V.五、；Bolmatov, D.(2014年3月)。“超临界氢的动态转变: 定义气体巨行星内部和大气之间的边界”。物理评论E。89(3) 032126。arXiv:1309.6500。Bibcode:2014PhRvE .. 89c2126T。doi:10.1103/PhysRevE.89.032126。PMID24730809。S2CID42559818。032126。  
\item 库尔特，Dauna。“木星内部的一种奇怪的流体？”。NASA。存档自原来的2021年12月9日。已检索12月8日，2021。
\item “NASA系统探索木星”。NASA。已存档从原始的2021年11月4日。已检索12月8日，2021。
\item 内疚，T。(1999)。“木星和土星内部的比较”。行星与空间科学。47(10-11):1183-1200。arXiv:astro-ph/9907402。Bibcode:1999年P & SS。47.1183克。doi:10.1016/S0032-0633(99)00043-4。S2CID19024073。已存档从原始的2021年5月19日。已检索6月21日，2023。 
\item 朗，肯尼斯·R。(2003)。木星: 一颗巨大的原始行星。NASA.存档自原来的2011年5月14日。已检索1月10日，2007。
\item 卡塔琳娜·洛德斯(2004)。“木星的焦油比冰多”(PDF)。天体物理学杂志。611(1):587-597。Bibcode:2004年4月。611 .. 587L。doi:10.1086/421970。S2CID59361587。存档自原来的(PDF)于2020年4月12日。  
\item Brygoo, S.;Loubeyre, P.;米洛，M.；Rygg, J.R、；Celliers，P。M、；Eggert, J.H、; Jeanloz，R.；柯林斯，G。W。(2021)。“木星内部条件下氢-氦不混溶的证据”。性质。593(7860):517-521.Bibcode:2021年自然593.517B。doi:10.1038/s41586-021-03516-0。OSTI1820549。PMID34040210。S2CID235217898。   
\item 克莱默，米里亚姆 (2013年10月9日)。钻石雨可能会填满木星和土星的天空。Space.com。已存档从原来的2017年8月27日。已检索8月27日，2017。
\item Kaplan，莎拉 (2017年8月25日)。“它在天王星和海王星上下雨固体钻石”。华盛顿邮报。已存档从原来的2017年8月27日。已检索8月27日，2017。
\item 吉洛，特里斯坦; 史蒂文森，大卫·J；威廉·B·哈伯德；索门，迪迪埃 (2004年)。“木星内部”。在弗兰的巴格纳尔; 道林，蒂莫西·E；威廉·B·麦金农 (主编)。木星.行星、卫星和磁层。剑桥行星科学。第1卷。英国剑桥: 剑桥大学出版社。第45页。Bibcode:2004jpsm.book...35G。ISBN 0-521-81808-7。已存档从原来的2023年3月26日。已检索3月19日，2023。
\item 阿特雷亚，苏西尔·K；Mahaffy，P。R、；尼曼，H. B.；王，M.H、; 欧文，T。C.(2003年2月)。“木星大气的组成和起源-更新，以及对太阳系外巨行星的影响”。行星与空间科学。51(2):105-112。Bibcode:2003P & SS...51 .. 105A。doi:10.1016/S0032-0633(02)00144-7。
\item 马克·J·洛夫勒；哈德森，雷吉L。(2018年3月)。“为木星的云着色: 氢硫化铵 (NH4SH) 的辐射分解”(PDF)。伊卡洛斯。302:418-425。Bibcode:2018年Icar .. 302 .. 418L。doi:10.1016/j.icarus.2017.10.041。已存档(PDF)从原来的2022年10月9日。已检索4月25日，2022年。
\item 英格索尔，安德鲁·P。；道林，蒂莫西·E；彼得·J·吉拉施；格伦·S·奥顿；阅读，彼得·L；桑切斯-拉维加，奥古斯丁; 表演者，亚当·P；艾米·西蒙·米勒；瓦萨瓦达，阿什温·R.(2004)。“木星的大气动力学”。在弗兰的巴格纳尔; 道林，蒂莫西·E；威廉·B·麦金农 (主编)。木星.行星、卫星和磁层(PDF)。剑桥行星科学。卷。1。英国剑桥: 剑桥大学出版社。pp。 105-128。Bibcode:2004jpsm.book .. 105I。ISBN  0-521-81808-7。已存档 (PDF)从原来的2022年10月9日。已检索3月8日，2022年。
\item 阿格利亚莫夫，尤里·S；鲁尼，乔纳森；贝克尔，海蒂N.；吉洛，特里斯坦; 吉伯德，塞兰G；阿特雷亚，苏西尔; 博尔顿，斯科特·J；史蒂文·莱文; 香农·T·布朗；黄，迈克尔·H (2021年2月)。“潮湿对流云中的闪电产生和对木星水量的限制”。地球物理研究杂志: 行星。126(2) e2020JE006504。arXiv:2101.12361。Bibcode:2021JGRE..12606504A。doi:10.1029/2020JE006504。S2CID231728590.E06504。 
\item 渡边，苏珊，主编 (2006年2月25日)。“令人惊讶的木星: 繁忙的伽利略航天器显示木星系统充满惊喜”。NASA.存档自原来的2011年10月8日。已检索2月20日，2007。
\item 克尔，理查德·A.(2000)。“深，湿热驱动木星天气”。科学。287(5455):946-947。doi:10.1126/science.287.5455.946b。S2CID129284864。已存档从原来的2023年2月3日。已检索4月26日，2022年。 
\item 贝克尔，海蒂N.；亚历山大，詹姆斯·W；阿特雷亚，苏西尔·K；斯科特·J·博尔顿；马丁·J·布伦南；布朗，香农T；亚历山大·纪尧姆; 特里斯坦·吉洛特; 安德鲁·P·英格索尔；史蒂文·M·莱文；鲁尼，乔纳森一世；阿格利亚莫夫，尤里·S；保罗·G·斯特菲斯.(2020)。“木星上的浅电风暴产生的小闪电”。性质。584(7819):55-58。Bibcode:2020年自然584. .. 55B。doi:10.1038/s41586-020-2532-1。ISSN0028-0836。PMID32760043。S2CID220980694。已存档从原始的2021年9月29日。已检索3月6日，2021。   
\item 吉洛，特里斯坦; 史蒂文森，大卫·J；阿特雷亚，苏西尔·K；斯科特·J·博尔顿；贝克尔，海蒂N.(2020)。风暴和木星中氨的消耗: I.“蘑菇” 的微观物理学”。地球物理研究杂志: 行星。125(8): e2020JE006403。arXiv:2012.14316。Bibcode:2020JGRE..12506403G。doi:10.1029/2020JE006404。S2CID 226194362。
\item 罗希尼·S·贾尔斯；Greathouse，托马斯·K；伯特兰·邦方; 格拉德斯通，G。兰德尔; 卡默，约书亚；Hue，文森特; Grodent，Denis C.；杰拉德，让-克劳德; Versteeg，Maarten H.; 黄，迈克尔H.; 博尔顿，斯科特J.；康纳尼，约翰E。P.;莱文，史蒂文M.(2020)。“在木星的高层大气中观察到的可能的瞬态发光事件”。地球物理研究杂志: 行星。125(11) e06659。arXiv:2010.13740。Bibcode:2020JGRE..12506659G。doi:10.1029/2020JE006659。S2CID225075904.E06659。 
\item Greicius，Tony，ed. (2020年10月27日)。“朱诺号的数据表明 '精灵' 或 '精灵' 在木星大气层中嬉戏”。NASA。已存档从原始的2021年1月27日。已检索12月30日，2020。
\item Strycker，P。D.;Chanover, N.;萨斯曼，M.；西蒙-米勒，A.(2006)。木星的发色团的光谱搜索。DPS第38号会议，#11.15。美国天文学会。Bibcode:2006DPS....38.1115S。
\item 彼得·J·吉拉施；尼科尔森，菲利普·D.(2004)。“木星”。世界图书 @ NASA。存档自原来的2005年1月5日。已检索8月10日，2006。
\item Chang，Kenneth (2017年12月13日)。“大红斑深入木星”。纽约时报。已存档从原来的2017年12月15日。已检索12月15日，2017。
\item 威廉·F·丹宁.(1899)。“木星，大红斑的早期历史”。皇家天文学会月刊。59(10):574-584。Bibcode:1899MNRAS..59..574D。doi:10.1093/mnras/59.10.574。
\item Kyrala, A.(1982)。“对木星的大红斑持续存在的解释”。月亮和行星。26(1):105-107.Bibcode:1982年M & P。26 .. 105k。doi:10.1007/BF00941374。S2CID121637752。 
\item 亨利·奥尔登堡主编 (1665-1666)。“皇家学会的哲学交易”。古腾堡计划。已存档从原来的2016年3月4日。已检索12月22日，2011年。
\item 黄，M.; 德佩特，我。(2008年5月22日)。“新的红点出现在木星上”。HubbleSite。NASA。已存档从原来的2013年12月16日。已检索12月12日，2013。
\item 西蒙·米勒，A.；Chanover, N.;奥顿，G.(2008年7月17日)。“三个红点混合在木星上”。HubbleSite。NASA。已存档从原来的2015年5月1日。已检索4月26日，2015。
\item 卡温顿，迈克尔A.(2002)。现代望远镜的天体。剑桥大学出版社.p. 53。ISBN 978-0-521-52419-3。
\item Cardall, C.Y.;Daunt, S.J.“大红斑”。田纳西大学。已存档从原来的2010年3月31日。已检索2月2日，2007。
\item 木星，太阳系的巨人。NASA.1979。第5页。已存档从原来的2023年3月26日。已检索3月19日，2023。
\item Sromovsky, L.A.;贝恩斯，K。H、; 弗莱，P。M、；卡尔森，R。W。(2017年7月)。“一种可能通用的红色发色团，用于模拟木星上的颜色变化”。伊卡洛斯。291:232-244。arXiv:1706.02779。Bibcode:2017Icar .. 291 .. 232S。doi:10.1016/j.icarus.2016.12.014。S2CID119036239。 
\item 怀特，格雷格 (2015年11月25日)。“木星的大红斑接近黄昏了吗？”。空间.新闻。已存档从原来的2017年4月14日。已检索4月13日，2017。
\item 索梅里亚，约尔; 迈耶斯，史蒂文·D；史威尼，哈里L.(1988年2月25日)。“木星的大红斑的实验室模拟”。性质。331(6158):689-693。Bibcode:1988 Natur.331..689S。doi:10.1038/331689a0。S2CID39201626。 
\item 艾米·西蒙；王，M.H、; 罗杰斯，J。H、; 奥顿，G。S、; 德佩特，I.；阿赛-戴维斯，X.；卡尔森，R。W.;马库斯，P。S。(2015年3月)。木星大红斑的戏剧性变化。第46届月球与行星科学会议。2015年3月16日至20日。德克萨斯州的伍德兰兹。Bibcode:2015LPI....46.1010S。
\item Grush，Loren (2021年10月28日)。NASA的朱诺宇宙飞船发现木星的大红斑有多深。The Verge。已存档从原始的2021年10月28日。已检索10月28日，2021。
\item 阿德里亚尼，阿尔贝托; 穆拉，A.；奥顿，G.；汉森，C.；Altieri，F.; 等人 (2018年3月)。“围绕木星的两极的旋风群”。性质。555(7695):216-219。Bibcode:2018 Natur.555..216A。doi:10.1038/自然25491。PMID29516997。S2CID4438233。  
\item 斯塔尔，米歇尔 (2017年12月13日)。“NASA刚刚看到木星上的大量旋风演变成令人着迷的六边形”。科学警报。已存档从原始的2021年5月26日。已检索5月26日，2021。
\item Steigerwald，Bill (2006年10月14日)。“木星的小红斑越来越强”。NASA.已存档从2012年4月5日起。已检索2月2日，2007。
\item 黄，迈克尔·H; 德佩特，Imke; 阿赛-戴维斯，西拉尔; 马库斯，菲利普·S；走吧，克里斯托弗·Y.(2011年9月)。“木星椭圆形BA变红前后的垂直结构: 发生了什么变化？”(PDF)。伊卡洛斯。215(1):211-225。Bibcode:2011伊卡。。215。211瓦。doi:10.1016/j.icarus.2011.06.032。已存档(PDF)从原来的2022年10月9日。已检索4月27日，2022年。
\item Stallard, Tom S.;梅林，亨里克; 米勒，史蒂夫; 摩尔，卢克; 奥多诺休，詹姆斯; 康纳尼，约翰E。P.;佐藤，武彦; 韦斯特，罗伯特·A；杰弗里·P·塞耶；徐，薇琪W.；约翰逊，罗西E.(2017年4月10日)。“木星的高层大气中的大冷点”。地球物理研究快报。44(7):3000-3008。Bibcode:2017乔治。。44.3000S。doi:10.1002/2016GL071956。PMC5439487。PMID28603321。  
\item 康纳尼，J。E.P.;Kotsiaros, S.;Oliversen, R.J.;埃斯普利，J.R、；乔根森，J。L.;乔根森，P。美国；Merayo, J.M。G.;赫尔塞格，M.；布洛瑟姆，J.；摩尔，K。M、；博尔顿，S.J.;莱文，S.M。(2017年5月26日)。“朱诺前九个轨道的木星磁场新模型”(PDF)。地球物理研究快报。45(6):2590-2596。Bibcode:2018乔治。。45.2590c。doi:10.1002/2018GL077312。已存档(PDF)从2022年10月9日开始。
\item 布雷纳德，吉姆 (2004年11月22日)。“木星的磁层”。天体物理学的旁观者。存档自原来的2021年1月25日。已检索8月10日，2008。
\item "无线电接收器"。NASA。2017年3月1日。存档自原来的2021年1月26日。已检索9月9日，2020。
\item 托尼·菲利普斯; 霍拉克，约翰·M。(2004年2月20日)。“木星上的无线电风暴”。NASA。存档自原来的2007年2月13日。已检索2月1日，2007。
\item Showalter, M.A.;伯恩斯，J.A.;Cuzzi, J.N.;波拉克，J.B.(1987)。“木星的环系统: 结构和粒子性质的新结果”。伊卡洛斯。69(3):458-498.Bibcode:1987年Icar。69 .. 458S。doi:10.1016/0019-1035(87)90018-2。
\item 伯恩斯，J.A.;Showalter, M.R、；汉密尔顿，D。P.;尼科尔森，P。D.; 德佩特，I.；奥克特-贝尔，M。E.;托马斯，P。C.(1999)。木星的微弱光环的形成科学。284(5417):1146-1150。Bibcode:1999Sci。284.1146B。doi:10.1126/science.284.5417.1146。PMID10325220。S2CID21272762。  
\item 菲斯勒，P。D.;亚当斯，O.W.;范德梅，N.；Theilig, E.E.;Schimmels, K.A.;刘易斯，G。D.;Ardalan, S.M、；亚历山大，C。J.(2004)。“伽利略星扫描仪在阿马尔西娅的观测”。伊卡洛斯。169(2):390-401。Bibcode:2004年Icar。。169。390F。doi:10.1016/j.icarus.2004.01.012。
\item 赫布斯特，T。M、；里克斯，H.-W。(1999)。“LBT上近红外干涉测量的恒星形成和太阳系外行星研究”。在盖特，艾克; 斯特克拉姆，布林弗里德; 克洛泽，西尔维奥 (编辑)。星际物质的光学和红外光谱。ASP会议系列。第188卷。加利福尼亚州旧金山: 太平洋天文学会。pp。 341-350。Bibcode:1999ASPC..188..341H。ISBN 978-1-58381-014-9。-见第3.4节。
\item 麦克道格尔，道格拉斯W.(2012年12月16日)。牛顿的引力: 宇宙力学的入门指南。施普林格纽约。第199页。ISBN 978-1-4614-5444-1。
\item 米奇琴科，T.A.;费拉兹-梅洛，S.(2001年2月)。“模拟木星-土星行星系统中的5:2平均运动共振”。伊卡洛斯。149(2):77-115.Bibcode:2001Icar .. 149 .. 357米。doi:10.1006/icar.2000.6539。
\item “模拟解释了具有偏心，近距离轨道的巨型系外行星”。科学日报。2019年10月30日。已存档从原来的2023年7月17日。已检索7月17日，2023。
\item "星际季节"。科学 @ NASA。存档自原来的2007年10月16日。已检索2月20日，2007。
\item Ridpath，伊恩(1998)。诺顿的星图(第19版)。Prentice Hall.ISBN 978-0-582-35655-9。[需要页面]
\item 藏起来，R.(1981年1月)。“关于木星的自转”。地球物理杂志。64:283-289。Bibcode:1981GeoJ。64。283小时。doi:10.1111/j.1365-246X.1981.tb02668.x。
\item 罗素，C。T.;于，Z。J.;Kivelson, M.G.(2001)。“木星的自转周期”(PDF)。地球物理研究快报。28(10):1911年-1912年。Bibcode:2001乔治。。28.1911r。doi:10.1029/2001GL012917。S2CID119706637。已存档(PDF)从原来的2022年10月9日。已检索4月28日，2022年。  
\item 罗杰斯，约翰·H (1995年7月20日)。“附录3”。巨行星木星。剑桥大学出版社。ISBN 978-0-521-41008-3。
\item 价格，弗雷德W。(2000年10月26日)。行星观察者手册。剑桥大学出版社。第140页。ISBN 978-0-521-78981-3。已存档从原来的2023年3月26日。已检索3月19日，2023。
\item 理查德·奥·菲梅尔；威廉·斯温德尔; 埃里克·伯吉斯 (1974)。“8.与巨人相遇”。先锋奥德赛(修订版)。NASA历史办公室。已存档从原来的2017年12月25日。已检索2月17日，2007。
\item 查普尔，格伦·F。(2009)。琼斯，劳伦诉；斯莱特，蒂莫西F. (编辑)。外行星。格林伍德引导宇宙。ABC-克里奥。第47页。ISBN 978-0-313-36571-3。已存档从原来的2023年3月26日。已检索3月19日，2023。
\item 北，克里斯; 亚伯，保罗 (2013年10月31日)。夜晚的天空: 如何解读太阳系。Ebury出版。第183页。ISBN 978-1-4481-4130-2。
\item 萨克斯，A。(1974年5月2日)。“巴比伦观测天文学”。伦敦皇家学会的哲学交易。276(1257): 43-50 (见第44页)。Bibcode:1974RSPTA.276...43S。doi:10.1098/rsta.1974.0008。JSTOR74273。S2CID121539390。  
\item 杜布斯，荷马H。(1958)。“中国天文学的开端”。美国东方学会杂志。78(4):295-300。doi:10.2307/595793。JSTOR595793。 
\item 陈，詹姆斯·L；陈，亚当 (2015)。哈勃太空望远镜物体指南: 它们的选择，位置和意义。施普林格国际出版社。第195页。ISBN 978-3-319-18872-0。已存档从原来的2023年3月26日。已检索3月19日，2023。
\item 大卫·西金特。J.(2010年9月24日)。“事实、谬论、不寻常的观察和其他杂项收集”。奇怪的天文学: 不寻常的故事，奇怪的，和其他难以解释的观察。天文学家的宇宙。pp. 221-282。ISBN 978-1-4419-6424-3。
\item Xi，Z。Z。(1981)。“甘德在伽利略之前2000年发现了木星的卫星”。天体物理学学报。1(2): 87。Bibcode:1981AcApS...1...85X。
\item 董，保罗 (2002)。中国的主要谜团: 超自然现象和中华人民共和国的未解之谜。中国图书。ISBN 978-0-8351-2676-2。
\item Ossendrijver，Mathieu (2016年1月29日)。“古巴比伦天文学家根据时间速度图下的区域计算出木星的位置”。科学。351(6272):482-484。Bibcode:2016Sci...351 .. 482O。doi:10.1126/科学.aad8085。PMID26823423。S2CID206644971。已存档从原来的2022年8月1日。已检索6月30日，2022年。  
\item Pedersen，Olaf (1974)。对Almagest的调查。欧登塞大学出版社，第423、428页。ISBN 978-87-7492-087-8。
\item Pasachoff, Jay M.(2015)。“西蒙·马吕斯的世界: 伽利略阴影下的400周年”。天文学史杂志。46(2):218-234。Bibcode:2015AAS...22521505P。doi:10.1177/0021828615585493。S2CID120470649。 
\item 威斯特法尔，理查德·S。“伽利略，伽利略”。伽利略计划。莱斯大学。已存档从最初的2022年1月23日。已检索1月10日，2007。
\item 德尔·桑托，保罗; 奥尔斯基，利奥·S。(2009)。“关于弗朗切斯科·丰塔纳给托斯卡纳·费迪南德二世·德·美第奇大公的一封未发表的信”。Galil æ ana: 伽利略研究杂志。六:1000-1017。已存档从原来的2023年11月15日。已检索11月14日，2023。 备用URL
\item 奥康纳，J。J.;罗伯逊，E。F.(2003年4月)。"乔瓦尼·多梅尼科·卡西尼"。圣安德鲁斯大学。已存档从原来的2015年7月7日。已检索2月14日，2007。
\item 大卫·H·阿特金森; 詹姆斯·B·波拉克；Seiff，Alvin (1998年9月)。“伽利略探测器多普勒风实验: 测量木星上的深纬向风”。地球物理研究杂志。103(E10):22911-22928。Bibcode:1998JGR...10322911A。doi:10.1029/98JE00060。
\item 保罗·默丁 (2000)。天文学和天体物理学百科全书。布里斯托: 物理研究所出版。ISBN 978-0-12-226690-4。
\item 罗杰斯，约翰·H (1995)。巨行星木星。剑桥大学出版社。pp. 188-189.ISBN 978-0-521-41008-3。已存档从原来的2023年3月26日。已检索3月19日，2023。
\item 理查德·奥·菲梅尔；威廉·斯温德尔; 埃里克·伯吉斯 (1974年8月)。“木星，太阳系的巨人”。先锋奥德赛(修订版)。NASA历史办公室。已存档从2006年8月23日起。已检索8月10日，2006。
\item 布朗，凯文 (2004)。“罗默假说”。MathPages。已存档从原来的2012年9月6日。已检索1月12日，2007。
\item Bobis，Laurence; Lequeux，James (2008年7月)。“卡西尼号、罗默号和光速”。天文历史与遗产杂志。11(2):97-105。Bibcode:2008JAHH...11...97B。doi:10.3724/SP.J.1440-2807.2008.02.02。S2CID115455540。 
\item Tenn，Joe (2006年3月10日)。“爱德华·爱默生·巴纳德”。索诺玛州立大学。存档自原来的2011年9月17日。已检索1月10日，2007。
\item "Amalthea概况介绍"。NASA/JPL。2001年10月1日。存档自原来的2001年11月24日。已检索2月21日，2007。
\item 小西奥多·邓纳姆 (1933)。“关于木星和土星光谱的注释”。太平洋天文学会的出版物。45(263):42-44。Bibcode:1933PASP...45...42D。doi:10.1086/124297。
\item 优素福，A.；马库斯，P。S。(2003)。“木星白色椭圆形从形成到合并的动态”。伊卡洛斯。162(1):74-93。Bibcode:2003年Icar。。162...74Y。doi:10.1016/S0019-1035(02)00060-X。
\item 温特劳布，瑞秋A.(2005年9月26日)。《田野中的一个夜晚如何改变了天文学》。NASA.存档自原来的2011年7月3日。已检索2月18日，2007。
\item 加西亚，伦纳德N.“木星十参无线电发射”。NASA.已存档从2012年3月2日起。已检索2月18日，2007。
\item 克莱因，M.J.;Gulkis, S.;博尔顿，S.J.(1996)。“木星的同步加速器辐射: 在SL9彗星撞击之前，之中和之后观察到的变化”。格拉茨大学会议。NASA: 217。Bibcode:1997版本4。conf .. 217k。已存档从原来的2015年11月17日。已检索2月18日，2007。
\item “先锋任务”。NASA.2007年3月26日。已存档从原来的2018年12月23日。已检索2月26日，2021。
\item “NASA格伦先锋发射历史”。NASA-格伦研究中心。2003年3月7日。存档自原来的2017年7月13日。已检索12月22日，2011年。
\item 彼得·W·福特斯库；史塔克，约翰; 斯温纳德，格雷厄姆 (2003)。航天器系统工程(第三版)。约翰·威利和儿子。第150页。ISBN 978-0-470-85102-9。
\item 平田，克里斯。“太阳系中的delta-v”。加州理工学院。存档自原来的2006年7月15日。已检索11月28日，2006。
\item Wong，Al (1998年5月28日)。"伽利略FAQ: 导航"。NASA.存档自原来的1997年1月5日。已检索11月28日，2006。
\item 陈，K.；帕雷德斯，E.美国；Ryne, M.S。(2004)。“尤利西斯的态度和轨道操作: 13年以上的国际合作”。2004年空间行动会议。美国航空航天学会。doi:10.2514/6.2004-650-447。
\item Lasher，劳伦斯 (2006年8月1日)。"先锋项目主页"。NASA太空项目部。存档自原来的2006年1月1日。已检索11月28日，2006。
\item “木星”。NASA/JPL。2003年1月14日已存档从原来的2012年6月28日。已检索11月28日，2006。
\item 汉森，C。J.;博尔顿，S.J.;Matson, D.L.;Spilker, L.J.;Lebreton，J.-P.(2004)。“木星的卡西尼-惠更斯飞越”。伊卡洛斯。172(1):1-8.Bibcode:2004年Icar。。172 .... 1小时。doi:10.1016/j.icarus.2004.06.018。
\item “冥王星的新视野看到了木星系统的变化”。NASA.2007年10月9日。存档自原来的2020年11月27日。已检索2月26日，2021。
\item “冥王星的新视野为木星系统提供了新的视角”。NASA.2007年5月1日。存档自原来的2010年12月12日。已检索7月27日，2007。
\item 麦康奈尔，香农 (2003年4月14日)。“伽利略: 木星的旅程”。NASA/JPL。存档自原来的2004年11月3日。已检索11月28日，2006。
\item 马加良斯，朱利奥 (1996年12月10日)。“伽利略探测任务事件”。NASA太空项目部。存档自原来的2007年1月2日。已检索2月2日，2007。
\item Goodeill，安东尼 (2008年3月31日)。"新边疆-任务-朱诺"。NASA.存档自原来的2007年2月3日。已检索1月2日，2007。
\item “朱诺，美国宇航局的木星探测器”。行星社会。已存档从2022年5月12日开始。已检索4月27日，2022年。
\item 喷气推进实验室 (2016年6月17日)。“美国宇航局的朱诺号宇宙飞船冒着木星的烟花为科学冒险”。phys.org。已存档从原来的2022年8月9日。已检索4月10日，2022年。
\item Firth，Niall (2016年9月5日)。美国宇航局的朱诺探测器拍摄了木星的北极的第一张照片。新科学家。已存档从原来的2016年9月6日。已检索9月5日，2016。
\item 克拉克，斯蒂芬 (2017年2月21日)。“美国宇航局的朱诺号航天器将保持在木星周围的当前轨道上”。现在太空飞行。已存档从原来的2017年2月26日。已检索4月26日，2017。
\item Agle, D.C.;温德尔，乔安娜; 施密德，Deb (2018年6月6日)。“NASA重新计划朱诺的木星任务”。NASA/JPL。已存档从原来的2020年7月24日。已检索1月5日，2019。
\item Talbert，Tricia (2021年1月8日)。NASA扩展了对两个行星科学任务的探索。NASA。已存档从原始的2021年1月11日。已检索1月11日，2021。
\item Dickinson，David (2017年2月21日)。“朱诺” 将停留在木星周围的当前轨道上。天空和望远镜。已存档从2018年1月8日的原件。已检索1月7日，2018。
\item Sori，Mike (2023年4月10日)。木星的卫星隐藏着巨大的地下海洋-两项任务正在发送航天器，看看这些卫星是否可以支持生命。对话。已存档从原来的2023年5月12日。已检索5月12日，2023。
\item 伯杰，布赖恩 (2005年2月7日)。“白宫缩减太空计划”。MSNBC。存档自原来的2013年10月29日。已检索1月2日，2007。
\item “拉普拉斯: 欧罗巴和木星系统的任务”。欧洲航天局。已存档从原来的2012年7月14日。已检索1月23日，2009年。
\item Fabio Favata (2011年4月19日)。“L级任务候选人的新方法”。欧洲航天局。已存档从原来的2013年4月2日。已检索5月2日，2012。
\item “欧洲航天局: 为木星冰冷的卫星任务发射升空”。BBC新闻。2023年4月14日已存档从原来的2023年4月14日。已检索4月14日，2023。
\item 杰夫·福斯特 (2020年7月10日)。“成本增长促使欧罗巴快船仪器发生变化”。太空新闻。已存档从原始的2021年9月29日。已检索7月10日，2020。
\item 琼斯，安德鲁 (2021年1月12日)。“中国的木星任务可能包括卡利斯特登陆”。行星社会。已存档从原始的2021年4月27日。已检索4月27日，2020。
\item 琼斯，安德鲁 (2021年4月16日)。“中国将向太阳系边缘发射一对航天器”。太空新闻。已存档从原始的2021年5月15日。已检索4月27日，2020。
\item 比林斯，李 (2019年11月12日)。“拟议的星际任务到达恒星，一次一代”。科学美国人。已存档从原始的2021年7月25日。已检索4月27日，2020。
\item 格林菲尔德博伊斯，内尔 (2023年2月9日)。“这就是为什么木星的月亮不断上升的原因”。NPR。已存档从原来的2023年3月5日。已检索3月29日，2023。
\item 谢泼德，斯科特·S.“木星的卫星”。卡内基科学研究所。已检索4月30日，2025。
\item 卡特，杰米 (2015)。初学者的观星程序。施普林格国际出版社。第104页。ISBN 978-3-319-22072-7。
\item Musotto, S.;瓦拉迪，F.；摩尔，W。B.;舒伯特，G.(2002)。“伽利略卫星轨道的数值模拟”。伊卡洛斯。159(2):500-504。Bibcode:2002年Icar。。159。500米。doi:10.1006/icar.2002.6939。存档自原来的2011年8月10日。已检索2月19日，2007。
\item 朗，肯尼斯·R。(2011年3月3日)。剑桥太阳系指南。剑桥大学出版社，第304页。ISBN 978-1-139-49417-5。
\item 露西·安·麦克法登; 保罗·魏斯曼; 托伦斯·约翰逊 (2006)。太阳系百科全书。爱塞唯尔科学。第446页。ISBN 978-0-08-047498-4。
\item 凯斯勒，唐纳德·J.(1981年10月)。“轨道物体之间碰撞概率的推导: 木星的外卫星的寿命”。伊卡洛斯。48(1):39-48。Bibcode:1981Icar...48...39K。doi:10.1016/0019-1035(81)90151-2。S2CID122395249。已存档从原始的2021年9月29日。已检索12月30日，2020。 
\item 汉密尔顿，托马斯W.M。(2013)。太阳系的卫星。SPBRA。第14页。ISBN 978-1-62516-175-8。
\item Jewitt, D.C.;谢泼德，S.；波科，C。(2004)。巴格纳尔，F.；道林，T.；麦金农，W. (编辑)。木星: 行星、卫星和磁层(PDF)。剑桥大学出版社。ISBN  978-0-521-81808-7。存档自原来的 (PDF)2009年3月26日
\item Nesvorný, D.;阿尔瓦雷洛，J。L.A.;多ones，L.；利维森，h.f.(2003)。“不规则卫星的轨道和碰撞演变”(PDF)。天文杂志。126(1):398-429。Bibcode:2003AJ....126..398N。doi:10.1086/375461。S2CID8502734。已存档(PDF)从原来的2020年8月1日。已检索8月25日，2019。  
\item "行星卫星平均轨道参数"。JPL,NASA。2013年8月23日已存档从原来的2013年11月3日。已检索2月1日，2016。,以及其中的引用。
\item Showman, A.P.;马尔霍特拉，R.(1999)。“伽利略卫星”。科学。286(5437):77-84。Bibcode:1999Sci...296...77S。doi:10.1126/science.286.5437.77。PMID10506564。S2CID9492520。  
\item 谢泼德，斯科特·S.;Jewitt, David C.(2003年5月)。“木星周围有大量不规则的小型卫星”(PDF)。性质。423(6937):261-263。Bibcode:2003 Natur.423..261S。doi:10.1038/nature01584。PMID12748634。S2CID4424447。存档自原来的(PDF)2006年8月13日   
\item 内斯沃尼，大卫; 博格，克里斯蒂安; 唐斯，卢克; 利维森，哈罗德·F。(2003年7月)。“不规则卫星家族的碰撞起源”(PDF)。天文杂志。127(3):1768-1783。Bibcode:2004AJ....127.1768N。doi:10.1086/382099。S2CID27293848。已存档(PDF)从2022年10月9日开始。  
\item 费拉兹-梅洛，S.(1994)。米兰尼，安德里亚; 迪·马蒂诺，米歇尔; 切利诺，A. (编辑)。柯克伍德间隙和共振群。小行星、彗星、流星1993年: 国际天文学联合会第160届专题讨论会论文集，1993年6月14日至18日在意大利贝尔吉拉特举行，国际天文学联合会。第160号专题讨论会。多德雷赫特: 克鲁沃学术出版社。第175页。Bibcode:1994IAUS..160..175F。
\item 克尔，理查德·A.(2004)。“木星和土星联队打击了内太阳系吗？”科学。306(5702): 1676。doi:10.1126/science.306.5702.1676a。PMID15576586。S2CID129180312。  
\item "木星木马列表"。IAU小行星中心。已存档从2011年7月25日起。已检索10月24日，2010。
\item 克鲁克申克，D。P.;Dalle Ore，C。M。；Geballe，T。R、；Roush, T.L.;欧文，T.C.;现金，米歇尔; 德伯格，C.；哈特曼，W。K。(2000年10月)。“特洛伊小行星624 Hektor: 表面组成的约束”。美国天文学会公报。32: 1027。Bibcode:2000DPS....32.1901C。
\item 奎因，T.；Tremaine, S.;邓肯，M。(1990)。“行星扰动和短周期彗星的起源”。天体物理学杂志，第1部分。355:667-679。Bibcode:1990apj。355 .. 667Q。doi:10.1086/168800。
\item “陷入困境: 火球照亮了木星”。科学日报。2010年9月10日。已存档从原来的2022年4月27日。已检索4月26日，2022年。
\item 中村，T.；Kurahashi，H. (1998)。“周期性彗星与地球行星的碰撞概率: 解析公式的无效情况”。天文杂志。115(2):848-854。Bibcode:1998AJ....115..848N。doi:10.1086/300206。
\item 霍纳，J.；琼斯，B。W。(2008)。“木星 -- 朋友还是敌人？I: 小行星 ”。国际天体生物学杂志。7(3-4):251-261.arXiv:0806.2795。Bibcode:2008IJAsB...7 .. 251小时。doi:10.1017/S1473550408004187。S2CID8870726。 
\item 再见，丹尼斯(2009年7月25日)。“木星: 我们的宇宙保护者？”。纽约时报。已存档从原来的2012年4月24日。已检索7月27日，2009年。
\item “深度 | P/鞋匠-征费9”。NASA太阳系探索。2017年12月9日.已存档从原来的2022年2月2日。已检索12月3日，2021。
\item 豪厄尔，伊丽莎白 (2018年1月24日)。“鞋匠-利维9号: 彗星的撞击在木星上留下了印记”。Space.com。已存档从原始的2021年12月6日。已检索12月3日，2021。
\item information@eso.org。“1994年在ESO进行的密集观测运动的大彗星坠毁”。eso.org。已存档从原始的2021年12月3日。已检索12月3日，2021。
\item “前20名彗星鞋匠-利维图像”。www2.jpl .nasa.gov。已存档从原始的2021年11月27日。已检索12月3日，2021。
\item 萨维奇，唐纳德; 艾略特，吉姆; 维拉德，雷 (2004年12月30日)。“哈勃观测揭示了木星碰撞的新线索”。nssdc.gsfc.nasa.gov。已存档从原始的2021年11月12日。已检索12月3日，2021。
\item “美国宇航局关于鞋匠-利维彗星的电视报道”。www2.jpl .nasa.gov。已存档从原始的2021年9月8日。已检索12月3日，2021。
\item Tabe，Isshi; Watanabe，jun-ichi; Jimbo，Michiwo (1997年2月)。“发现1690年记录的对木星可能的撞击点”。日本天文学会的出版物。49:L1-L5。Bibcode:1997PASJ...49L...1T。doi:10.1093/pasj/49.1.l1。
\item “观星者为木星日光视图做准备”。ABC新闻。2005年6月16日。存档自原来的2011年5月12日。已检索2月28日，2008。
\item 罗杰斯，J。H. (1998年)。“古代星座的起源: I.美索不达米亚传统 ”。英国天文协会杂志。108:9-28.Bibcode:1998JBAA..108....9R。
\item 韦尔登，B。L.(1974)。“古巴比伦天文学”(PDF)。科学觉醒II。多德雷赫特: 斯普林格。pp. 46-59。doi:10.1007/978-94-017-2952-9_3。ISBN  978-90-481-8247-3。已存档 (PDF)从原来的2022年3月21日。已检索3月21日，2022年。
\item “行星的希腊名字”。2010年4月25日已存档从原来的2010年5月9日。已检索7月14日，2012。在希腊语中，木星的名字是Dias，这是神宙斯的希腊名字。 另请参见希腊关于地球的文章。
\item 西塞罗，马库斯·图利乌斯 (1888)。西塞罗的Tusculan争议; 此外，关于神的性质和英联邦的论文。由Yonge，查尔斯·杜克翻译。纽约，纽约: 哈珀兄弟公司。 274-通过互联网档案。
\item 马库斯·图卢斯·西塞罗(1967) [1933]。沃明顿，E。H、 (编)。De Natura Deorum[论神的本质]。西塞罗.19卷。翻译: 拉克姆，H.剑桥，马萨诸塞州: 剑桥大学出版社。 175-通过互联网档案。
\item 佐洛特尼科娃，O.(2019年)。“接触中的神话: 荷马宙斯中的syro-phoenician特征”。科学遗产。41(5):16-24.已存档从原来的2022年8月9日。已检索4月26日，2022年。
\item Tarnas, R.(2009)。“行星”。Archai: 原型宇宙学杂志。1(1):36-49.CiteSeerX10.1.1.456.5030。 
\item 哈珀，道格拉斯 (2001年11月)。“木星”。在线词源词典。已存档从原来的2008年9月28日。已检索2月23日，2007。
\item Vytautas tu ė nas (2016)。“波罗的海雷神Perkunas与他的古董等价物Jupiter和Zeus之间的共同属性”(PDF)。地中海考古学和考古学。16(4):359-367。已存档(PDF)从原来的2023年7月19日。已检索7月19日，2023。
\item "Guru"。印度n Divinity.com。已存档从原来的2008年9月16日。已检索2月14日，2007。
\item Sanathana, Y.美国；曼吉尔，哈扎里卡 (2020年11月27日)。“雅鲁藏布江山谷的占星术: 对纳瓦格拉哈雕塑描绘的反思”(PDF)。遗产: 考古学多学科研究杂志。8(2):157-174.已存档(PDF)从原来的2022年10月9日。已检索7月4日，2022年。 [死链接]
\item "图尔克Astrolojisi-2"(土耳其语)。NTV。存档自原来的2013年1月4日。已检索4月23日，2010。
\item Jan Jakob Maria De Groot (1912)。中国的宗教: 大学主义。道家与儒学研究的关键。美国关于宗教史的讲座。第10卷。G.P.普特南的儿子。第300页。已存档从原来的2024年2月26日。已检索1月8日，2010。
\item 克鲁普，托马斯 (1992)。日本的数字游戏: 现代日本对数字的使用和理解。日产研究所/Routledge日本研究系列。Routledge. pp. 39-40.ISBN 978-0-415-05609-0。
\item Hulbert，Homer Bezaleel (1909)。韩国的通过。Doubleday, Page & Company. p. 426。已检索1月8日，2010。
\end{enumerate}
\subsection{进一步阅读}
\begin{itemize}
\item Bagenal, Fran；Dowling, Timothy E.；McKinnon, William B.（2006）。《木星：行星、卫星与磁层》（Jupiter: The Planet, Satellites and Magnetosphere）。剑桥大学出版社（Cambridge University Press）。ISBN 978-0-52-103545-3 —— 经由 Google Books。
\item Lohninger, Hans；等（2005 年 11 月 2 日）。“由旅行者 1 号所见的木星”（"Jupiter, As Seen By Voyager 1"）。收录于《太空之旅》（A Trip into Space）。虚拟应用科学研究所（Virtual Institute of Applied Science）。
\end{itemize}
\subsection{外部链接}
\begin{itemize}
\item NASA 科学任务理事会（Science Mission Directorate）发布的木星概览（Jupiter overview by NASA's Science Mission Directorate）
\item 《NASA 的朱诺号探测器如何改变我们对木星的全部认知》（How NASA's Juno probe changed everything we know about Jupiter）。《Scientific American》，2025 年 8 月 19 日。
\item Tony Dunn 的 Gravity Simulator 网站上的木星 62 颗卫星的模拟（Simulation of the 62 moons of Jupiter from Tony Dunn's Gravity Simulator website）。
\item 加州大学圣克鲁兹分校图书馆数字馆藏中利克天文台档案的 1920 年代木星照片（Photographs of Jupiter circa 1920s from the Lick Observatory Records Digital Archive, at the University of California, Santa Cruz Library's Digital Collections）。
\item 木星系统的交互式三维引力模拟（Interactive 3D gravity simulation of the Jovian system）。
\item 业余天文学家拍摄的 2010 年木星撞击事件的 YouTube 视频（Video on YouTube of the 2010 Jupiter impact event by an amateur astronomer）。
\item “Jupiter”——BBC Radio 4《In Our Time》节目（"Jupiter" – In Our Time, BBC Radio 4）
\item NASA 制作的朱诺号飞掠木卫三及木星的 YouTube 动画（Animation on YouTube of the Juno spacecraft's flyby of Ganymede and Jupiter by NASA）。
\end{itemize}